% ============================================================================
%  DisloCluster Code History
%
%  A numbered record of every change made from MoDELib2-NNL onward, beginning
%  at the point where the immobile population equations were introduced into
%  it to make it consistent with the ZrMicro 0-D model, and running through to
%  the present state of the repository.
%
%  Build:  pdflatex "DisloCluster Code History.tex"   (three times, for the
%          table of contents and the cross references)
%  Self-contained: no bibliography, no external figures.
% ============================================================================
\documentclass[11pt,a4paper]{article}

\usepackage[T1]{fontenc}
\usepackage[utf8]{inputenc}
\usepackage{lmodern}
\usepackage[margin=2.3cm]{geometry}
\usepackage{amsmath,amssymb}
\usepackage{booktabs}
\usepackage{array}
\usepackage{longtable}
\usepackage{xcolor}
\usepackage[colorlinks=true,linkcolor=blue!50!black,urlcolor=blue!50!black,
            citecolor=blue!50!black]{hyperref}

% Long \code{} identifiers cannot be hyphenated, so give the paragraph builder
% room to absorb them rather than letting them run into the margin.
\setlength{\emergencystretch}{3em}

\newcommand{\code}[1]{\texttt{\small #1}}
\newcommand{\dpa}{\ensuremath{\mathrm{dpa}}}
\newcommand{\aloop}{\ensuremath{\langle a\rangle}}
\newcommand{\cloop}{\ensuremath{\langle c\rangle}}
\newcommand{\Om}{\ensuremath{\Omega}}

% ---- the numbered-change environment --------------------------------------
% A local stand-in for needspace, which is not in every MiKTeX install: break
% the page if fewer than #1 lines remain, so a change heading is never orphaned
% from its first paragraph. On an empty page \pagegoal is \maxdimen and the
% test correctly does nothing.
\newcommand{\needlines}[1]{%
  \par\dimen0=\pagegoal \advance\dimen0 by -\pagetotal
  \ifdim\dimen0<#1\baselineskip \newpage \fi}

\newcounter{chg}
\newcommand{\change}[1]{%
  \needlines{4}%
  \refstepcounter{chg}%
  \par\medskip\noindent
  {\bfseries\boldmath Change \thechg.\quad #1}%
  \par\nobreak\smallskip}
\newcommand{\why}[1]{\par\noindent\textbf{\itshape Why.}\ #1\par}
\newcommand{\impact}[1]{\par\noindent\textbf{\itshape Code.}\ #1\par}
\newcommand{\cost}[1]{\par\noindent\textbf{\itshape Cost.}\ #1\par}
\newcommand{\status}[1]{\par\noindent\textbf{\itshape Status.}\ #1\par}

\title{\bfseries DisloCluster Code History\\[2mm]
       \large Every change from MoDELib2-NNL to the present, numbered,\\
       beginning with the immobile population equations}
\author{DisloCluster development note}
\date{15 August 2026}

\begin{document}
% This document is dense with long, unhyphenatable \code{} identifiers.
% Prefer a slightly loose line to one that runs into the margin.
\sloppy
\maketitle

\begin{abstract}
\noindent
This document records, in one continuous numbering, every substantive change
made on the way from \code{MoDELib2-NNL} to DisloCluster. Each entry states
\emph{what} was changed, \emph{why}, what it cost or bought in \emph{code}
structure, and what it did to \emph{computational efficiency}.

The baseline is the point at which the \textbf{immobile population equations
were introduced} into MoDELib2-NNL --- which until then had no immobile
species at all --- to make the 3-D code consistent with the ZrMicro 0-D model:
loop nucleation, absorption growth, vacancy-loop annealing, loop--loop
coalescence and loop--network coalescence. Those occupy
Part~\ref{part:immobile} and are Changes~\ref{chg:baseline}--\ref{chg:burgers}.
Everything after them follows from consequences that only became visible once
the immobile families were alive: the mobile solve had to be repaired to feed
them (Part~\ref{part:mobile}), the material file had to be made to agree with
the 0-D calibration (Part~\ref{part:material}), a genuine field-level coupling
had to replace a four-scalar handoff (Part~\ref{part:coupling}), the
conservation bookkeeping had to be reinterpreted for an open system
(Part~\ref{part:physics}), and the cost of the resulting scheme had to be
brought down by a factor of several (Part~\ref{part:performance}).

Two conventions are used throughout. First, a change is reported with the
number that was \emph{measured}, not the number that was expected; where a
change bought nothing, that is stated. Second, correctness checks that were
available in principle but not exercised are not claimed as verification.
The open items are collected, unresolved, in Part~\ref{part:open}.
\end{abstract}

\tableofcontents
\newpage

% ===========================================================================
\section{The baseline: MoDELib2-NNL, before the immobile populations}
% ===========================================================================

\textbf{The baseline for this document is
\code{mlm335/MoDELib2-NNL}}, the state of the code at the moment the immobile
population equations were introduced into it. That is where the work recorded
here begins, and Change~\ref{chg:baseline} is that introduction.

At that point the code solved the mobile steady-state cluster-dynamics PDE and
\emph{nothing else}. Its immobile species count \code{iSize} was \textbf{0},
\code{solveImmobileClusters()} was an empty function body, there was no
\code{ImmobileSinkRate.h} and no immobile integrator of any kind, and the
source carried a literal \code{// Missing immobile sinks} placeholder at
\code{ClusterDynamicsFEM.cpp:110} where the loop absorption term belongs. The
loop populations were not present in a simplified form --- they were
\textbf{absent}. Loops appeared in the model only as a hole where a sink term
should be.

The immobile equations below were therefore written into that code to make it
consistent with the \textbf{ZrMicro 0-D model}: cascade nucleation, absorption
growth, vacancy-loop annealing, and coalescence in the loop--loop and
loop--network channels, each given the same meaning it carries in ZrMicro so
that the 3-D solve reduces to the 0-D result in the well-mixed limit. The
governing equations are those of Po \& Ghoniem, Deliverable D1/M1, \S2.2;
equation numbers below refer to that document.

\paragraph{Above the baseline.} This checkout records that the tree descends
from \textbf{Po's MoDELib2} --- \code{README.md} reads \code{\# MoDELib2} and
\code{Contributors.txt} is headed by Giacomo Po. The registered submodule URL
was \code{https://github.com/Ghoniem/MoDELib2-NNL.git}. Neither upstream is
vendored here, so nothing further about the ancestry can be verified from the
repository itself.

\paragraph{A note on \code{mlm335/MoDELib-fullCD}, and on an earlier error in
these documents.} \code{MoDELib-fullCD} is a \emph{separate} cluster-dynamics
branch: \code{iSize} is 8, and it carries \code{ImmobileSinkRate.h},
\code{SpatialODESolver.h}, \code{FirstOrderReaction.h} and a
continuum\,$\rightarrow$\,discrete loop conversion. Comparing the two
repositories after the fact shows 31 cluster-dynamics parameter keys shared
verbatim and a set of systematic renames (\code{loopNucDefects} $\rightarrow$
\code{nNuc}, \code{loopCoalLL/LN} $\rightarrow$ \code{cLL/cLN},
\code{loopCoalKappa*} $\rightarrow$ \code{kappa*}, \code{loopAnnealTau0\_SI}
$\rightarrow$ \code{tau0\_vLoop\_SI}, \code{loopCoalNetwork\_SI} $\rightarrow$
\code{rhoNetwork\_SI}, \code{minimumLoopSize} $\rightarrow$ \code{r\_min}).

An earlier revision of these notes read that overlap as \emph{descent} and
declared fullCD the correct baseline, describing the immobile solver below as
a re-discretization of fullCD's. \textbf{That inference was wrong.} Shared
naming across two branches of one code, developed in one group from one
formulation, is evidence of \emph{convergence}, not of derivation --- and it
says nothing about which came first. The immobile equations here were written
to reproduce ZrMicro, not to re-express fullCD.

fullCD remains worth comparing against, and this document does so wherever the
comparison is informative --- most usefully in
Change~\ref{chg:rediscretize}, where its Galerkin/mass-matrix discretization
is set beside the nodal one adopted here, and in Part~\ref{part:open}, where
its continuum\,$\rightarrow$\,discrete loop conversion is noted as a
capability this branch does not have. Those are comparisons between siblings,
not a record of what was inherited.

\medskip\noindent
\textbf{Chronology.} Changes~\ref{chg:baseline}--\ref{chg:usecase} were made
on the \code{zr3d\_ghoniem} branch of the fork before the repository was
consolidated. Changes~\ref{chg:fieldchannel} onward are dated by the
repository's own history, from \code{c7574f6} (7~August 2026) to
\code{6dbcc68} (15~August 2026).

% ===========================================================================
\part{The immobile population equations}
\label{part:immobile}
% ===========================================================================

\textbf{This is where the work started, and it started from nothing.}
MoDELib2-NNL had no immobile species: the loop populations that the ZrMicro
0-D model is calibrated on had no representation in the 3-D code at all. Every
kinetic channel in this part --- cascade nucleation, absorption growth,
vacancy-loop annealing, loop--loop coalescence, loop--network coalescence ---
was written into the code here, given the same meaning it carries in ZrMicro,
and then integrated by a scheme that survives the stiffness those channels
bring with them.

The changes divide into two kinds, and it is worth separating them because
they failed in completely different ways.
Changes~\ref{chg:baseline}--\ref{chg:coalescence} \emph{introduce} physics: a
channel that was not there before is added and given its 0-D form.
Changes~\ref{chg:gate}, \ref{chg:anneal-implicit}, \ref{chg:anneal-content},
\ref{chg:emission} and \ref{chg:coal-implicit} then \emph{repair how those
channels were integrated} --- a term written in a form that could not enter an
implicit update, a unit conversion missing so a gate never closed, or a
channel integrated explicitly when its own time constant sat orders of
magnitude below the step.

The second kind produced by far the larger errors. A missing channel is
usually visible as a trend that is simply wrong; a channel present but
integrated explicitly past its stability limit produces a population that
resets itself every step, which looks like physics until the numbers are
checked against an analytic balance. The loop content was wrong by a factor of
350 for that reason, not because the annealing term was absent.

% ---------------------------------------------------------------------------
\section{Introducing the immobile populations}
% ---------------------------------------------------------------------------

\change{\code{iSize} 0 $\rightarrow$ 8: the immobile families are created}
\label{chg:baseline}
The immobile degrees of freedom are introduced as four families $\times$ two
quantities each --- number density and defect content:
$[N_c,\,N_{a1},\,N_{a2},\,N_{a3}\;|\;c_c,\,c_{a1},\,c_{a2},\,c_{a3}]$.
One vacancy-type \cloop\ family and three prismatic interstitial-type \aloop\
families.
\why{The 0-D model carries exactly these populations, and the point of the
exercise is that the 3-D reduce to it in the well-mixed limit. This is the
change that brings the immobile half of the model into existence: before it,
\code{iSize} was 0 and there was no loop population to evolve, no sink for the
mobile species to see, and nothing for the ZrMicro parameter set to attach
to.}
\impact{\code{ClusterDynamicsParameters.h}. The change is global --- because
\code{iSize} is a compile-time constant threaded through every
cluster-dynamics template, \code{Library/Materials/Zr4.txt} can no longer be
run from this branch at all. That is why the Zr3d\_ghoniem material definition
is a \emph{new} file rather than an edit of \code{Zr4.txt}, which is left
untouched.}
\cost{The immobile DOF vector is $8$ per node against $4$ mobile, so the
stored state per node roughly triples. The immobile solve itself is nodal and
matrix-free, so this is memory, not time.}

\change{The immobile time integration scheme}
\label{chg:rediscretize}
With the families created, an integrator for them had to be written. The
scheme adopted is nodal collocation with sub-cycled, fused semi-implicit
updates. The sibling \code{MoDELib-fullCD} branch made different choices at
each of the three decision points, and is set beside it here because the
contrast is what makes the reasoning legible --- not because one was derived
from the other:

\begin{center}
\begin{tabular}{lll}
\toprule
 & \code{MoDELib-fullCD} (sibling) & this branch \\
\midrule
Rate & Galerkin: assembled at quadrature points, then & nodal collocation; \\
     & $L^2$-projected through the consistent mass    & no mass matrix \\
     & matrix (CG to $10^{-4}$, \code{SpatialODESolver}) & \\[1mm]
Time update & explicit Euler \code{dof += rate*dt}, then & one \emph{fused}
                semi-implicit \\
            & \emph{sequential} implicit loss factors  & update
                $\dfrac{n+\Delta t\,\dot n_{\rm nuc}}{1+\Delta t\,\lambda}$ \\
            & $(1+\Delta t\lambda_k)^{-1}$ per channel & \\[1mm]
Sub-cycling & none --- one step of \code{dtMax} & \code{nSub = 20} per dose
              step \\
\bottomrule
\end{tabular}
\end{center}

\why{Three reasons, in order of weight. (i) The loop equations at a point
depend only on that point's own state and the local mobile field --- there is
no spatial operator in them --- so the mass-matrix projection buys consistency
of a quantity that has no spatial coupling to be consistent about. (ii)
Sequential per-channel loss factors do not commute, so the answer depends on
the order the channels are applied; a fused update does not. (iii) One step of
\code{dtMax} is not enough: at $G=10^{-7}$~\dpa/s a dose step is
$\sim7\times10^{19}$ MoDELib time units, and several channels have time
constants far below that.}
\impact{\code{solveImmobileClusters()} in
\code{src/ClusterDynamics/ClusterDynamicsFEM.cpp} --- an \emph{empty function
body} in MoDELib2-NNL --- is written as a node loop with an inner sub-step
loop. DOF layout is \code{index = node*iSize +
component}, with $[0,n_F)$ densities and $[n_F,2n_F)$ contents. No mass
matrix, no linear solve, no assembly: the immobile step is matrix-free.}
\cost{This is the choice that makes the whole later coupling affordable.
Collocating at nodes avoids a conjugate-gradient solve per dose step
altogether; the 20 sub-steps cost 20 evaluations of an algebraic update at
each node, which is arithmetic with no memory traffic beyond the node's own
$8$ DOFs. Measured later against the CVODE march, this nodal integrator is
first-order semi-implicit and \emph{has} to be: the measured stiffness puts
its step $46\,000$ times beyond the forward-Euler stability limit. That
measurement also corrected a description of it, common in earlier notes, as
explicit Euler --- it is not, and could not be.}

% ---------------------------------------------------------------------------
\section{Nucleation}
% ---------------------------------------------------------------------------

\change{Cascade-borne loop nucleation, Eqs.\ (36)/(51)}
\label{chg:cascade}
Loops are seeded from the displacement cascade at
$\dot n_{\rm nuc} = G_0\,\varepsilon_k / n_{\rm nuc}(k) / \Om$, with the
matching content source $G_k = G_0\varepsilon_k$, where $\varepsilon_k$ is the
family's share of the cascade (\code{loopCascadeFractions}) and $n_{\rm nuc}$
is the number of defects a newborn loop carries (\code{nNuc}).
\why{This is the 0-D's \code{nucleation\_rate\_iL}/\code{aiL}, and it must be
reproduced term for term if the two codes are to agree.}
\impact{New parameters \code{loopCascadeFractions} and \code{nNuc} in
\code{ClusterDynamicsParameters}, each guarded by \code{iSize>0} and converted
to MoDELib units on read.}
\cost{Negligible --- two multiplications per family per sub-step.}

\change{Homogeneous SIA clustering as a second nucleation source}
\label{chg:clustering}
\aloop\ loops are also nucleated by interstitial clustering reactions whose
product falls off the end of the mobile ladder. The channels are selected
automatically by a new \code{getLoopNucChannels()}, which takes the pairs
$(a,b)$ from the existing \code{reactionMap} that (i) have reactants of the
same polarity, and (ii) have a product size $m_a+m_b$ carried by \emph{no}
mobile species. For $\{v,i,2i,3i\}$ that selects $i{+}3i$, $2i{+}2i$ and
$2i{+}3i$, and rejects $i{+}i \rightarrow 2i$, $i{+}2i \rightarrow 3i$ and
every vacancy-bearing pair.
\why{\aloop\ loop density measured $0.53\times$ the 0-D at 30~\dpa\ and flat in
dose, while the \cloop\ density matched to 0.1\% at every dose --- so the
defect was specific to the interstitial families. The 0-D nucleates \aloop\
loops from \emph{two} sources, and only the cascade one was wired:

\begin{center}
\begin{tabular}{lrr}
\toprule
source & loops/atom/s & share \\
\midrule
cascade $G\varepsilon_{iL}/n_{iL,{\rm nuc}}$ & $7.0429\times10^{-13}$ & 61.6\% \\
clustering $R_{i,3i}+R_{2i,2i}$              & $4.3969\times10^{-13}$ & 38.4\% \\
\bottomrule
\end{tabular}
\end{center}

\noindent
Cascade-only nucleation therefore predicts a density ratio of 0.62 against the
0.53 measured --- the right size and sign for the entire discrepancy. The
interstitials those reactions consume were simply disappearing: with no $4i$
species to receive the product, \code{getR2()} debits the reactants and credits
nothing. That is the same truncation that raises
\code{Warning: Sum of R2 is not zero} at start-up, which had until then been
recorded as benign. It is not benign; it is a nucleation source.}
\impact{New member \code{loopNucChannels} and method
\code{getLoopNucChannels()}. Taking the flux from the reaction network the
mobile solve \emph{already uses} makes the split mass-conserving by
construction: the loops gain exactly what the mobile field loses. Two traps
had to be navigated. (i) \code{getR2()} writes $-K_{ab}$ into \emph{both}
\code{R2[a](a,b)} and \code{R2[a](b,a)}, so $c^{T}R_2 c = -2K_{ab}c_ac_b$ ---
but the residual is assembled as \code{R2*(0.5*mobileClusters)}, the
quadratic-form factor that makes \code{bWF\_R2 = R2*c} its exact Jacobian, so
the per-species loss rate is $K_{ab}c_ac_b$ and not twice that. (ii)
Cascade-borne loops are born at \code{nNuc} defects while clustering-borne
embryos carry only the $|m_a|+|m_b|$ atoms of the event that made them, so the
two sources must \emph{not} be divided by the same number. The family split is
in proportion to \code{loopCascadeFractions} within each polarity, which
reproduces the 0-D's own $f_{na}/f_a$ split exactly and generalizes to any
number of families.}
\cost{One pass over a three-entry channel list per node per dose step, hoisted
out of the sub-step loop because the mobile field is frozen during it.
Immeasurable.}
\status{Verified against the 0-D jump frequencies
$\omega_m = z_c\nu\exp(-E_m/kT)$ at 573~K:

\begin{center}
\begin{tabular}{lrrr}
\toprule
channel & Zr3d\_ghoniem $K$ [1/s] & ZrMicro & error \\
\midrule
$i+3i \rightarrow 4i$  & $5.0535398\times10^{7}$ & $\omega_i+\omega_{3i}=5.0535398\times10^{7}$ & $4\times10^{-6}$ \\
$2i+2i \rightarrow 4i$ & $9.8765170\times10^{2}$ & $4\omega_{2i}=9.8812394\times10^{2}$ & 0.05\% \\
$2i+3i \rightarrow 5i$ & $4.9382578\times10^{2}$ & $\omega_{2i}+\omega_{3i}=4.9406197\times10^{2}$ & 0.05\% \\
\bottomrule
\end{tabular}
\end{center}

\noindent
Effect at 30~\dpa, median interior: \aloop\ loop density $0.533 \rightarrow
\mathbf{0.704}$, content $0.620 \rightarrow \mathbf{0.784}$, defects per loop
$1.164 \rightarrow \mathbf{1.114}$, all as ratios to the 0-D. The \cloop\
density was unaffected at $0.999$, as it must be.}

% ---------------------------------------------------------------------------
\section{Growth by absorption}
% ---------------------------------------------------------------------------

\change{Absorption growth with the diffusional-anisotropy bias, Eqs.\ (41)--(45)}
\label{chg:absorption}
Each family grows by absorbing mobile defects of like polarity and shrinks by
absorbing the opposite kind:
$\dot g_k = \phi^{\rm like}_k - \text{gate}\cdot\phi^{\rm opp}_k$, with
$\phi_k = \sum_m \bar D_m Z(\text{row}(k),m)\,S_k\,c_m\,|m_m|$ --- the
size weight $|m_m|$ being what makes a $2i$ cluster deliver two interstitials.
\why{This is the growth flux of the 0-D's \code{loop\_growth\_rate\_*}
methods.}
\impact{New method \code{loopDADbias()} returning the Eq.~(15) capture
efficiencies $Z_c = Z^0 p$, $Z_a = Z^0(p+p^{-2})/2$, row~0 for \cloop\ and
row~1 for \aloop, from two new fitted parameters \code{dadAnisotropy} ($p_m$)
and \code{dadZ0} ($Z^0_m$). The orientationally averaged diffusivity
$\bar D_m = (\det D_m)^{1/3}$ is formed once per solve.}
\cost{$n_F\times m_{\rm size}$ multiply--adds per node per sub-step; the
dominant arithmetic in the immobile step, and still trivially cheap.}
\status{The DAD reconciliation is a \emph{fitting} form, not a generated one
--- see Part~\ref{part:open}.}

\change{The minimum-stable-size shrink gate, Eq.\ (91)}
\label{chg:gate}
The shrinking channel is switched off smoothly for loops below the minimum
stable radius: a cubic Hermite gate in $x = (r_k/r_{\min}-1)/0.3$, zero at
$x\le0$ and one at $x\ge1$.
\why{Below $r_{\min}$ a loop is not a stable sink; letting it shrink further
is unphysical and drives the radius to zero.}
\impact{Reads the existing \code{r\_min} key --- but \emph{that key was never
converted from metres}. It was read raw in SI and compared against radii from
\code{rloop()}/\code{rpyr()}, which are in units of $b$. The gate was
therefore permanently wide open and the change had no effect until
\code{r\_min} was divided by \code{b\_SI} on read. The conversion is safe
because \code{r\_min} is used only by the new code.}
\cost{One branch and a cubic per family per sub-step.}

\change{Loop densities stored per $b^3$, not per atom}
\label{chg:perb3}
The immobile number-density DOF is stored as loops per $b^3$; the content DOF
is a defect volume fraction.
\why{\code{clusterRadius()}, \code{clusterDensity()} and \code{sigmoid()} all
form $n = c/(N\Om)$ --- i.e.\ they \emph{expect} $N$ in loops per $b^3$. The
new integration was accumulating loops per atom.}
\impact{\code{nucRate} is divided by \Om, \code{Nvol} is $n(k)$ directly, and
the like-loop spacing is $d_{LL} = (1/N)^{1/3}$. Content needs no conversion:
a volume fraction is numerically equal to the 0-D per-atom content.}
\cost{None. But the consequence of getting it wrong was not cosmetic: the
helpers read the mean loop size as $2.82\times$ too large, which was enough to
push the \cloop\ family across the bi-pyramid\,$\rightarrow$\,loop sigmoid onto
the wrong branch. A units error in a quantity that feeds a branch condition
changes which physics runs, not just by how much.}

% ---------------------------------------------------------------------------
\section{Annealing and thermal emission}
% ---------------------------------------------------------------------------

\change{Vacancy-loop dissolution integrated implicitly}
\label{chg:anneal-implicit}
The $n/\tau_{vL}$ and $c/\tau_{vL}$ dissolution channels (Eqs.~38/53) were
moved from the explicit source term into the implicit denominator.
\why{\cloop\ loops were collapsing onto the positivity floor: $n_{vL}$ pinned
at $10^{-30}$ while \aloop\ loops nucleated and grew normally. The cause was
integration, not physics --- $\tau_{vL}\approx9.95\times10^{16}$ MoDELib time
units against a sub-step of $3.48\times10^{18}$, so \textbf{the decay time is
35 times shorter than the step}. Explicit Euler gives $n(1-35)<0$, clamped
every iteration.

The confirmation that this was integration and not a modeling error is the
analytic balance: $n^* = \dot n_{\rm nuc}\tau_{vL} = 1.79\times10^{-8}$,
against the 0-D's $C_{vL}+C_{avL} = 1.78\times10^{-8}$. The physics was
already right.}
\impact{The linear losses are collected into a single coefficient
\code{lossN} and applied as $(n + \Delta t\,\dot n_{\rm nuc})/(1+\Delta t\,
\text{lossN})$.}
\cost{One division instead of one multiplication. The scheme becomes
unconditionally stable in this channel and returns $n^* = \dot n_{\rm nuc}
\tau_{vL}$ exactly in the limit $\Delta t/\tau \rightarrow \infty$ --- so the
sub-step count is set by accuracy in the \emph{other} channels rather than by
stability in this one. Verified: $n_{vL} = 1.79161\times10^{-8}$.}

\change{Annealing releases each loop's \emph{birth} content, not its mean}
\label{chg:anneal-content}
The content sink from vacancy-loop dissolution was changed from $c/\tau_{vL}$
(proportional to content) to $n_{\rm nuc}\,n\,\Om/\tau_{vL}$ (proportional to
\emph{density}).
\why{\cloop\ loop content came out 350 times below the 0-D, with a mean size
$m_c = 487$ against $n_{\rm nuc} = 400$ --- i.e.\ pinned at the nucleation
size --- versus $1.7\times10^{5}$ in the 0-D. ZrMicro's
\code{annealing\_content\_vL()} returns $n_{vL,{\rm nuc}}\,C_{vL}/\tau_{vL}$:
each dissolving loop releases the content it was \emph{born} with, because
$\tau_{vL}$ describes the dissolution of still-small embryo loops, while loops
that survive and grow are stable. Removing the \emph{mean} content instead
caps the mean size at $\sim n_{\rm nuc}$, which is exactly what was observed.}
\impact{A separate \code{contentSink} term, explicit (it is proportional to
$n$, not to $c$, so it cannot enter the $c$-implicit denominator), added to the
content source with a negative sign.}
\cost{Arithmetically identical. The self-check is the valuable part: at number
saturation $n^* = \dot n_{\rm nuc}\tau_{vL}$, so the release equals
$n_{\rm nuc}\dot n_{\rm nuc}\Om = G_k$ and \emph{exactly cancels the cascade
content seed}, leaving net content growth flux-driven --- precisely what the
0-D docstring describes. A term that is right in form can be checked against
an identity like that; one that is right only in magnitude cannot.}
\status{\cloop\ content moved from $350\times$ low to within 20--26\% of the
0-D.}

\change{Thermal emission written as a coefficient, not a rate}
\label{chg:emission}
The Eq.~(55)--(56) single-vacancy emission term was recast from an absolute
rate into an equivalent first-order loss coefficient.
\why{As an absolute rate it could not enter the implicit update at all, which
left the stiffest available channel explicit.}
\impact{Per loop, $\alpha_v = 2\pi r D_v/\Om \cdot e^{-E_b/kT}$ (the
convention of \code{getR1()}); the content lost per unit volume is
$\alpha_v N \Om = S_k D_v e^{-E_b/kT}$, the \Om's cancelling. Since that does
not scale with $c$, it is divided by $c$ to form an equivalent first-order
coefficient and added to \code{lossC}.}
\cost{One extra division per vacancy family per sub-step. It buys the ability
to integrate the term implicitly, which is the point.}

% ---------------------------------------------------------------------------
\section{Coalescence}
% ---------------------------------------------------------------------------

\change{Geometric coalescence in two channels, Eqs.\ (97)--(102)}
\label{chg:coalescence}
Two Avrami-type coalescence channels were implemented:

\begin{align*}
\text{loop--loop:}\quad &
  \phi_{LL} = 1-\exp\!\left(-\kappa_{LL}\tfrac{4}{3}\pi r_{LL}^3 N\right),
  \qquad \nu_{LL} = c_{LL}\,v_{\rm abs}\,N^{1/3} \\
\text{loop--network:}\quad &
  \phi_{LN} = 1-\exp\!\left(-\kappa_{LN}\pi r_{LN}^2 \rho_N\right),
  \qquad \nu_{LN} = c_{LN}\,v_{\rm abs}\,\sqrt{\rho_N}
\end{align*}

\noindent
with $r_{LL} = \min(r_k, N^{-1/3})$, $r_{LN} = \min(r_k, \rho_N^{-1/2})$, and
the climb speed taken from the \emph{absorbed} flux,
$v_{\rm abs} = \phi^{\rm like}_k / (2\pi r_k N)$.
\why{These are the 0-D's coalescence channels, and they are the dominant loop
\emph{removal} mechanism --- without them the loop density grows without bound
in dose.}
\impact{New parameters \code{cLL}, \code{cLN}, \code{kappaLL},
\code{kappaLN}, \code{rhoNetwork}. \textbf{The two channels are not
symmetric}, and the asymmetry is physical: the like-loop channel merges two
loops into one, so it removes number but \emph{conserves content}; the
loop--network channel absorbs a loop into the dislocation network, so it
removes \emph{both}. Hence \code{coalN = nuLL*phiLL + nuLN*phiLN} but
\code{coalC = nuLN*phiLN} alone.

Taking the climb speed from the absorbed flux rather than the net one matters
at saturation, where the net flux passes through zero while coalescence is
still active; $v_{\rm abs}$ stays positive there.}
\cost{Two exponentials, a cube root and a square root per family per sub-step
--- the most expensive arithmetic in the immobile update, and still far below
the cost of a single sparse matrix--vector product.}

\change{Coalescence moved into the implicit update}
\label{chg:coal-implicit}
Both coalescence channels were reclassified as first-order loss
\emph{coefficients} and moved into the implicit denominator.
\why{The same stiffness failure as Change~\ref{chg:anneal-implicit}, in the
channel that had been left explicit. \code{nuLN*phiLN*dt\_sub} reaches
$\sim3.6$, so each sub-step overshoots and clamps to the floor. The signature
was unmistakable: \aloop\ mean loop size $m_a = 5042$ at 1~\dpa, then
\emph{exactly} $300.0$ --- $= n_{\rm nuc}$ --- at 2~\dpa. Content collapsing
$11\times$ while density kept climbing is not a physical trajectory; it is a
population being reset to its nucleation state every step.}
\impact{Eq.~(97) removes loops in proportion to $n$ and content in proportion
to $c$, so both are already first-order coefficients and the move is exact,
not an approximation. They are added to \code{lossN} and \code{lossC}
respectively.}
\cost{No arithmetic change. \aloop\ loop density went from 3--4$\times$ the
0-D to within $\sim$30\%.}

\change{\aloop\ Burgers vectors given as lattice components}
\label{chg:burgers}
\code{immobileSpeciesBurgers} was corrected from Cartesian to lattice
components: $(1,0,0)$, $(0,1,0)$, $(-1,1,0)$ for the three \aloop\ families
and $(0,0,1)$ for \cloop.
\why{At \emph{zero} applied stress the three prismatic families must be
identical by symmetry. They measured $n_{a1},n_{a2},n_{a3} = 2.84\times
10^{-7}$, $2.32\times10^{-7}$, $3.16\times10^{-7}$.}
\impact{The columns of \code{immobileSpeciesBurgers} are multiplied by
\code{latticeBasis} on read --- they are lattice components, and Cartesian
ones had been supplied. The resulting magnitudes were $|b| = 1.0$, $0.753$,
$1.197$ instead of all $1.0$. After the fix,
\code{immobileSpeciesBurgersMagnitude: 1.633, 1.0, 1.0, 1.0}.}
\cost{None.}
\status{This established a \textbf{standing test}: at zero applied stress the
three \aloop\ family densities must agree. It has been checked on every run
since.}

% ===========================================================================
\part{The mobile solve, and the coupling between the two}
\label{part:mobile}
% ===========================================================================

The immobile families do not evolve in isolation: they absorb from the mobile
field, and the mobile field must see them as sinks. Half the changes in this
part exist because turning on the immobile side exposed defects in the mobile
side that had been invisible while the immobile populations were zero.

\change{\code{ImmobileSinks.h} --- the sink the mobile species see}
\label{chg:immobilesinks}
A new \code{EvalFunction}, mirroring the existing \code{SecondOrderReaction},
supplying the loop absorption sink that the source had marked
\code{// Missing immobile sinks}. For each mobile species $m$ it assembles
\[
  k^2_m = \sum_k Z(\text{type}(k), m)\, S_k , \qquad
  S_k \ \text{from}\ \code{clusterDensity()} = 2\pi r_k N_k
\]
and returns $-k^2_m \bar D_m$ on the diagonal, matching the sign convention of
\code{getR1()}.
\why{Without it the coupling is one-way: loops grow by absorbing mobile
defects that are never removed from the mobile pool. That is not a small
error at dose --- it is a violated mass balance that grows with the loop
population.}
\impact{New file \code{include/ClusterDynamics/ImmobileSinks.h}, wired into the
mobile Newton system in \code{ClusterDynamicsFEM.cpp}: \code{bWF\_RI} into the
Jacobian and \code{lWF\_RI} into the residual, alongside the existing
\code{R1}/\code{R2} terms.

Two sinks are deliberately \emph{absent} from this file. The grain boundary is
imposed spatially by the Dirichlet boundary condition on the mobile species and
must not be double-counted here. The network dislocation sink reuses the
existing \code{otherSinks\_SI} plumbing, which required no new code at all.}
\cost{One more bilinear weak form in the Newton assembly --- of the same order
as \code{R1}. This is a real per-iteration cost on the fast solve, and it is
one of the terms later hoisted out of the Newton loop
(Change~\ref{chg:hoist}).}

\change{\code{loopSinkScale} --- reproducing the 0-D sink convention exactly}
\label{chg:loopsinkscale}
An optional material key, default $(1,1,1,1)$, multiplying the loop sink
strength: $\code{loopSinkScale} = (0.291528,\ 0.792317,\ 0.792317,\ 0.792317)$.
\why{ZrMicro's loop absorption is \emph{not} the purely geometric
$S_k = 2\pi r_k N_k$, and it departs from it in two independent ways.

\emph{(a) One radius prefactor for all four families.} \code{input\_data.py}
derives both $l_a = \sqrt{\Om/\pi b_a}$ and $l_c = \sqrt{\Om/\pi b_c}$, but
\code{reaction\_rates.loop\_absorption()} uses $l_c/l$ for \emph{every}
family --- $l_a$ is computed and never used anywhere in the codebase. The 3-D
builds $r_k$ from each family's own $|b_k|$, so each needs
$\sqrt{|b_k|_{\rm MoDELib}/b_c^{\rm ZrMicro}}$: for \cloop,
$\sqrt{5.279467/5.15} = 1.012495$; for \aloop, $\sqrt{3.233/5.15} = 0.792317$.

\emph{(b) A fitted $Q$ on the \cloop\ channel only:}
$1.012495\times0.287931 = 0.291528$.}
\impact{Applied \emph{identically} in \code{ImmobileSinks.h} (the sink seen by
the mobile species) and in \code{solveImmobileClusters()} (the growth flux it
causes) --- mirroring the 0-D, where \code{loop\_absorption()} and the
\code{loop\_growth\_rate\_*} methods share one prefactor. Applying it in only
one place would break the mass-conserving coupling between them, which is
worth stating because it is the kind of asymmetry that produces a slow
conservation drift rather than an obvious failure.}
\cost{One multiplication.}

\change{\code{Eb} declared before \code{R1} --- an inherited upstream bug}
\label{chg:eb}
The declaration of \code{Eb} was moved ahead of \code{R1} in the class.
\why{$C_{2i}$ and $C_{3i}$ came out at $\sim10^{-17}$ against the 0-D's
$\sim2.5\times10^{-11}$ --- six orders low. \code{getR1()} reads \code{Eb(k)}
to build the cluster dissociation rates, but \code{Eb} was declared
\emph{after} \code{R1}. C++ initializes members in \textbf{declaration} order,
not in the order of the constructor's initializer list, so this was an
uninitialized read. It evaluated to $\exp(0)=1$ instead of $\exp(-E_b/kT)$,
inflating the $2i$/$3i$ dissociation rate to the bare attempt frequency.

The evidence was arithmetic rather than circumstantial: \code{R1(2,2)} $=
-5.0535\times10^{7}$~1/s where the network sink alone gives $-20.86$;
\code{R1(3,3)/R1(2,2)} $= 0.4802 = p(1,2)/p(1,1)$ exactly, so \emph{both}
dissociation rates had lost their Boltzmann factor; and the predicted
$C_{2i}$ under the bug, $G_{2i}/\alpha = 9.89\times10^{-18}$, matched the
observed $1.06\times10^{-17}$.}
\impact{A declaration move, with the reasoning recorded at the declaration
site so it cannot be undone by a later tidy-up.}
\cost{None. \code{R1(2,2)} $\rightarrow -21.05$~1/s; $C_{2i}$ $+1.2\%$ and
$C_{3i}$ $-2.3\%$ against the 0-D.}
\status{\textbf{This is a bug in the original code, not in the Zr3d\_ghoniem
changes.} It was invisible in Po's calibration because $E_b = 5.0$~eV for
every species makes $\exp(-5/kT) \sim 10^{-44}$, negligible whether or not the
factor is applied. Giving $2i$ its physical $0.957$~eV is what exposed it. It
affects any MoDELib run with a non-negligible cluster binding energy and is
worth reporting upstream.}

\change{Clusters transport at the monomer diffusivity}
\label{chg:clusterdiff}
$E_m$ for $2i$ and $3i$ was set to the monomer value, and the
$\omega_{2i}/\omega_i$ factor moved into the reaction prefactors instead.
\why{Setting $E_m^{2i} = 1.36292$~eV drops $D_{2i}$ to $\sim5\times10^{-6}$ of
$D_i$, and \code{FixedDirichletSolver failed.} was thrown from
\code{compute()} --- the factorization, not the solve. The mobile solve is a
\emph{pure steady-state} (QSSA) system: \code{dmBWF} is diffusion only, with no
mass term, so the $2i$/$3i$ block had nothing on its diagonal but a vanishing
Laplacian. The matrix was singular.

The fix is also \emph{more} faithful to the 0-D, which is the part worth
recording. ZrMicro is inconsistent about which frequency it applies where, but
it is unambiguous: \code{R\_2i\_s}/\code{R\_3i\_s} (network sinks) and
\code{flux\_i} (loop absorption) all use $\omega_i$, and $\omega_{2i}$ appears
\emph{only} in the mobile--mobile cluster reactions.}
\impact{Material file change plus the reaction-prefactor derivation of
Change~\ref{chg:prefactors}. The network-sink and loop-absorption terms now
reproduce the 0-D exactly, which they did not before.}
\cost{Removes a singular block, which is the difference between a solve and no
solve.}

\change{The reaction prefactors derived and verified}
\label{chg:prefactors}
All ten entries of \code{reactionPrefactorMap} recomputed from
\[
  p_{ab} = f_{ab}\,\frac{z_c}{a^2}\,
           \frac{\Om}{4\pi(r_a+r_b)}\,
           \frac{D^{0D}_a+D^{0D}_b}{D^{3D}_a+D^{3D}_b}
\]
\why{MoDELib forms $K_{ab} = p_{ab}4\pi(r_a+r_b)(D_a+D_b)/\Om$ in
\code{getR2()}, while ZrMicro forms $K_{ab} = f_{ab}(\omega_a+\omega_b)$.
Equating them is the only way the two reaction networks can agree. The last
ratio is unity except for pairs involving $2i$/$3i$, where it restores the
$\omega_{2i}$-based rate now that those species transport at the monomer
diffusivity (Change~\ref{chg:clusterdiff}).}
\impact{Material file only.}
\cost{None.}
\status{\textbf{Confirmed empirically.} MoDELib prints the assembled
second-order interaction matrices in 1/s at start-up:

\begin{center}
\begin{tabular}{lrrr}
\toprule
entry & MoDELib & ZrMicro & agreement \\
\midrule
$(v,i)$  & $2.526876\times10^{7}$  & $0.5(\omega_i+\omega_v) = 2.52560\times10^{7}$ & 0.05\% \\
$(i,i)$  & $-2.021161\times10^{8}$ & $-4\omega_i = -2.020216\times10^{8}$ & 0.05\% \\
$(v,2i)$ & $6.940722\times10^{3}$  & $\omega_{2i}+\omega_v = 6.93473\times10^{3}$ & 0.09\% \\
\bottomrule
\end{tabular}
\end{center}

\noindent
This also settles a factor-of-two ambiguity that could not be resolved by
reading the code: \code{SecondOrderReaction.h} documents the rate as
$\dot c = \frac12 R_2 cc$ while \code{getR2()} fills both $(a,b)$ and $(b,a)$,
and it was not clear whether the $\frac12$ survives into the weak form. The
match above shows it does.}

\change{\code{atomicVolume\_SI} --- an atomic volume independent of the cell volume}
\label{chg:atomicvol}
A new optional material key, read by a new static
\code{getClusterAtomicVolume()}.
\why{\code{Polycrystal::Omega} is $\det(\text{latticeBasis}) = 1.41421\,b^3$
for HEX --- the volume of a cell containing \textbf{two} atoms, not the atomic
volume the cluster-dynamics equations assume. Every
content~$\leftrightarrow$~radius~$\leftrightarrow$~density conversion inherited
the error.}
\impact{Made \emph{optional} deliberately: material files that omit the key
keep the previous behavior, so \code{Polycrystal} --- which the elastic and
dislocation-dynamics sides also use --- is left alone. This is the one place
where a global correction would have reached beyond cluster dynamics, and it
was scoped rather than applied globally.}
\cost{None directly, but see Change~\ref{chg:omegacancel}, which this caused,
and Change~\ref{chg:omega}, which corrected the \emph{value} much later.}

\change{\code{mBWF}/\code{dmBWF} scale with \code{cdp.omega}, not \code{poly.Omega}}
\label{chg:omegacancel}
\why{Introduced by Change~\ref{chg:atomicvol}. \code{FluxMatrix} returns
$-D/\code{cdp.omega}$, and the weak forms multiplied it back by
\code{ddBase.poly.Omega}. Those two were the same number until the CD atomic
volume became independently settable; afterwards the diffusion term carried a
spurious factor of $\code{poly.Omega}/\code{cdp.omega} = 1.41421/0.355105 =
\mathbf{3.98}$ against every reaction term.}
\impact{Use \code{cdp.omega} in both weak forms, restoring the exact
cancellation so the flux is $-D\nabla c$.}
\cost{None. Found by reading the constructor while chasing the
non-convergence of Change~\ref{chg:clamp} --- a factor of four on the
diffusion operator is a plausible contributor to that, and this is a reminder
that a correction in one place (an atomic volume) can silently unbalance a
cancellation somewhere else entirely.}

\change{The mobile fixed-point iteration: five solver controls, and a diagnosis}
\label{chg:clamp}
Five optional \code{DD.txt} scalars were added to the mobile fixed-point loop,
all defaulting to the historical behavior except the iteration cap:
\code{mobileSolverTolerance} ($10^{-5}$), \code{mobileSolverMaxIterations}
(200; \code{0} = unlimited, the old behavior), \code{mobileSolverRelaxation}
(1.0), \code{mobileSolverErrorMode} (0), \code{mobileSolverClampInLoop} (1).
\why{\code{solveMobileClusters} iterates \code{while(cError>cTol)} with
\textbf{no iteration cap and no stagnation test}. That code is inherited
verbatim from MoDELib2-NNL, and the sibling \code{MoDELib-fullCD} carries the
identical loop, so it is not a defect introduced here. It is safe only while
the iteration always converges. Solving the mobile
field against an immobile state supplied by an external march, \code{cError}
falls from $5.8\times10^{5}$ to $1.3\times10^{-4}$ in nine iterations and then
\emph{stagnates}, oscillating between $4.7\times10^{-5}$ and
$1.2\times10^{-4}$. With no cap, DDomp never returns.

The cause is a clamp added \emph{here}, not anything inherited. Neither
MoDELib2-NNL nor fullCD has a mobile clamp at all; this branch inserted
\code{clampMobileClusters()} between the Newton update and the error test,
which turns the iteration into a \emph{projected} Newton method.

The precise failure is worth stating, because the obvious explanation is
wrong. The active set does \emph{not} oscillate: an instrumented run shows it
freeze at 7090 DOFs by iteration 3 and never change again. What fails is that
the Newton matrix is the \emph{unconstrained} Jacobian --- it does not know any
DOF is pinned. Each iteration it computes an increment for all DOFs including
the pinned ones; the projection then discards exactly those moves, but the
increment for the \emph{free} DOFs was computed assuming the pinned ones would
move. The free DOFs respond every iteration to a coupling that is immediately
cancelled, and that inconsistency does not decay. Solve-then-project is simply
not a consistent method for a constrained system.}
\impact{Each iteration now prints the iteration number, which species attains
the maximum, and how many DOFs the floor moved. On hitting the cap the
\emph{best} iterate seen is restored, not the last.}
\cost{Measured from an identical saved immobile state (24\,115 CD nodes,
96\,460 mobile DOFs):

\begin{center}
\small
\begin{tabular}{@{}lrlrr@{}}
\toprule
configuration & iters & converged & final \code{cError} & wall \\
\midrule
clamp in loop, undamped (as shipped) & 30 (capped) & no  & $6.4\times10^{-5}$, cycling & 392~s \\
\textbf{clamp deferred out of the loop} (fullCD) & \textbf{5} & \textbf{yes} & $\mathbf{6.83\times10^{-7}}$ & \textbf{112~s} \\
clamp in loop, damped $w=0.5$ & 30 (capped) & not yet & $9.4\times10^{-5}$ & 400~s \\
clamp in loop, damped $w=0.7$ & 20 & yes & $8.2\times10^{-6}$ & 283~s \\
clamp in loop, \code{errorMode=1} & 40 (capped) & no & $5.1\times10^{-5}$ & 492~s \\
clamp in loop, \code{cTol=1e-4} & 19 & ``yes'' & $9.2\times10^{-5}$ & 258~s \\
\bottomrule
\end{tabular}
\end{center}

\noindent
Deferring the clamp restores \textbf{quadratic} convergence ---
$2.95\times10^{4} \rightarrow 1.23\times10^{3} \rightarrow 6.77\times10^{1}
\rightarrow 5.73\times10^{-2} \rightarrow 6.83\times10^{-7}$ --- i.e.\ a
genuine Newton method, and it is faster than the solve has ever been:
$3.5\times$ against the shipped configuration.

The accuracy consequence is larger than the speed one. Maximum relative
difference in $C_i$ against the converged field, over the 22\,344 nodes above
the floor: damped $w=0.7$, \textbf{0.19\%}; \code{errorMode=1}, 21.1\%;
\code{cTol=1e-4}, 49.9\%; baseline, \textbf{53.6\%}. Two configurations that
\emph{report} success are wrong by half. Relaxing the tolerance is not a fix;
it declares victory on the bad field.}
\status{Both routes must use \code{mobileSolverClampInLoop=0} --- the coupled
march sets it in \code{modelib\_qssa}, and the standalone comparison driver in
\code{setup\_standalone.\allowbreak FAST\_STEP\_SETTINGS}. Left at the default on one side
only, a route-to-route comparison measures this solver defect rather than the
physics.}

\change{\code{useImmobileSolver} --- running the fast step alone}
\label{chg:useimmobile}
A new \code{DD.txt} scalar, default 1 when the key is absent so every
pre-existing case is unaffected, making the \code{solveImmobileClusters()}
call conditional.
\why{\code{ClusterDynamicsFEM::solve()} already \emph{is} the two-time-scale
split: \code{solveMobileClusters()} is the fast step (steady $C_M^*(x)$, no
time derivative anywhere) and \code{solveImmobileClusters()} is the slow one.
To let an external driver own the slow step --- which is what the whole
coupled march is --- the second call must be suppressible.}
\impact{With \code{useImmobileSolver=0}, DDomp performs the fast step and
passes the immobile field through untouched. Because \code{runSingleStep}
writes output \emph{before} incrementing \code{runID}, a run with
\code{startAtTimeStep=0} and \code{Nsteps=1} reads \code{evl\_0.txt} and
overwrites the same file --- an in-place round trip, which is exactly the
interface a Python driver wants.

\code{modelib\_qssa.py} asserts on the \code{immobile solver SKIPPED} line, so
a binary predating this flag fails \emph{loudly} rather than silently advancing
the immobile field twice. That assertion is the only thing standing between a
stale binary and a run that looks fine and is wrong by a factor of two in
dose.}
\cost{This is the enabling change for the entire coupled architecture of
Part~\ref{part:coupling}.}

% ===========================================================================
\part{The material file and the verification case}
\label{part:material}
% ===========================================================================

\change{\code{Library/Materials/Zr3d\_ghoniem.txt} --- a new material definition}
\label{chg:material}
Everything outside the \code{\# ClusterDynamics} section (elasticity, mobility
laws, gamma surface) is inherited unchanged; original values are preserved in
comments in the file itself. The changed values:

\begin{center}
\small
\begin{tabular}{@{}p{4.1cm}p{2.2cm}p{3.0cm}p{5.2cm}@{}}
\toprule
Key & Zr4 (Po) & Zr3d\_ghoniem (ZrMicro) & Basis \\
\midrule
\code{atomicVolume\_SI} & \emph{absent} & \code{1.2e-29} & ZrMicro \code{physical\_props\allowbreak['Omega']} \\
\code{mobileSpecies\allowbreak SurvivingEfficiency} & 0.01 & 1.0 & the 0-D loses defects only to the in-cascade cluster fractions \\
\code{mobileSpecies\allowbreak CascadeFractions} & \code{1.0 0.6 0.3212 0.0788} & \code{0.995 0.97788714 0.005 0.00333333} & $G_v/G = 1-\varepsilon_{vL}$, $G_i/G = 1-\varepsilon_{2i}-\varepsilon_{3i}-\varepsilon_{iL}$, etc. \\
\code{mobileSpecies\allowbreak EnergyFormation\_eV} & \code{1.5} & \code{1.6} & 0-D $E_F^v$; gives $C_v^{\rm eq} = 8.45\times10^{-15}$ at 573~K \\
\code{mobileSpecies\allowbreak EnergyMigration\_eV} & \code{0.9 / 0.09 / …} & \code{1.2 / 0.759101 (×3)}, isotropic & 0-D $E_m^v$, $E_m^i$; clusters at the monomer value (Change~\ref{chg:clusterdiff}) \\
\code{mobileSpecies\allowbreak D0\_SI} & \code{4.9639e-10}, \code{8.1882e-10} & \code{5.21645e-7} all, isotropic & $D_0 = a^2\nu$, $a = 3.23\times10^{-10}$~m, $\nu = 5\times10^{12}$~s$^{-1}$ \\
\code{reactionPrefactorMap} & \code{2.856}, \code{1.0}, \code{0.0} & all ten recomputed & Change~\ref{chg:prefactors} \\
\code{Eb\_eV} & \code{5.0 5.0 5.0 5.0} & \code{5.0 5.0 0.957033 2.0} & index 0 $\rightarrow$ loop emission; 2, 3 $\rightarrow$ 0-D $E_b^{2i}$, $E_b^{3i}$ \\
\code{otherSinks\_SI} & \code{0 0 0 0} & \code{1.88e14}, \code{1.8988e14 ×3} & network dislocation sink, $k^2_v = \rho_N$, $k^2_i = Z_N\rho_N$ \\
\bottomrule
\end{tabular}
\end{center}

\why{Every one of these is a term the 0-D carries; leaving any of them at the
Po calibration means the two codes are solving different models and no
comparison between them means anything.}
\impact{A new \S2.2 block adds \code{immobileSpeciesVector},
\code{immobileSpeciesRelRelaxVol}, \code{immobileSpeciesBurgers},
\code{alpha\_bp}, \code{delVPyramid}, \code{w0}, \code{n\_s}, \code{evc},
\code{Nvmax}, \code{nmin}, \code{nmax}, \code{r\_min}, \code{distanceFactor},
\code{loopCascadeFractions}, \code{nNuc}, \code{cLL}, \code{cLN},
\code{kappaLL}, \code{kappaLN}, \code{tau0\_vLoop\_SI}, \code{Ea\_vLoop\_eV},
\code{rhoNetwork\_SI}, \code{dadAnisotropy}, \code{dadZ0}. The coalescence
coefficients, nucleation sizes and dissolution constants are the fitted 0-D
values from run \code{20260622\_144021\_7959445}.

The migration-energy row order matters and is easy to get wrong: components
are $[11\;12\;13\;22\;23\;33]$, and the off-diagonal $D_0$ entries are zero, so
the tensors are diagonal.}
\cost{None.}
\status{\code{Zr4.txt} is \textbf{untouched} and carries no immobile keys ---
but that is a property of this one file and \textbf{not} evidence about any
other branch. fullCD's \code{Zr4\_Fitted.txt} and \code{Zr4\_BMD19\_nuc.txt}
each carry a full immobile key set. \code{Zr4.txt}'s missing keys have been
cited as evidence about upstream in earlier notes; they are not.}

\change{\code{tutorials/zrmicro\_coupled/} --- the verification case}
\label{chg:usecase}
A cluster-dynamics-only verification case: \code{useFEM=1},
\code{useDislocations=0}, \code{useClusterDynamics=1},
\code{dtMax = 6.95868964e19} (= 1~\dpa\ at $G = 10^{-7}$),
\code{outputFrequency=1}, \code{startAtTimeStep=0},
\code{absoluteTemperature=573}, non-periodic, zero target loop densities.
\why{A fixed, reproducible case is the only way to attribute a change in the
answer to a change in the code.}
\impact{\code{clean\_run.sh} refreshes the material file and regenerates
\code{evl\_0} on every invocation. Four separate traps made that necessary,
and each is worth recording because each failed \emph{silently}:

\begin{enumerate}\itemsep2pt
\item \emph{\code{microstructureGenerator} never terminates.} It is the
  discrete-dislocation path: it inserts loops one at a time, and the log showed
  it had reached $2.34\times10^{17}$~m$^{-3}$ against a target of
  $1.48\times10^{21}$ --- 0.016\% of the way, needing $\sim10^{11}$ discrete
  loops. Cluster-dynamics loop densities are not representable as discrete
  dislocations, which is precisely why they belong in the continuum fields.
  Fixed with \code{useDislocations=0} and \code{targetDensity=0}, giving an
  empty 20-byte \code{evl\_0.txt} in about a second. DDomp still requires that
  file to exist even with dislocations off.
\item \emph{Restart silently double-counted dose.} \aloop\ loop content came
  out at exactly $3\times$ the expected value. \code{startAtTimeStep=-1} means
  ``last available step'', and \code{evl\_0.txt} serves double duty as both the
  initial configuration \emph{and} the runID~0 output, so a bare re-run resumed
  from the previous run's accumulated dose. Diagnosed by
  $c_a^{\max} = 2.112\times10^{-3}$ = exactly $3\times(G_0\varepsilon_a\cdot
  1\,\dpa)$.
\item \emph{Stale material file.} \code{inputFiles/} holds a \emph{copy} of the
  material file, made before the \S2.2 block existed; the run died on
  \code{does not cointain line with format immobileSpeciesVector=…}.
\item \emph{CRLF line endings produced a zero-node mesh.} \code{unitCube\_15K.msh}
  checked out with CRLF and the \code{.msh} parser returned zero nodes
  \emph{without complaint}, surfacing much later as a periodicity failure.
\end{enumerate}}
\cost{None, but the recurring hazard is worth stating once: \textbf{almost
every failure in this stack exits 0} --- the build script, DDomp with a missing
\code{evl\_0}, a missing material key, a CRLF mesh, and the solver failures
above. Verification has been by artifacts on disk and log tails, never by exit
status. Later, \code{build.preflight()} was written to \emph{run} each binary
rather than trust its presence, for the same reason.}

% ===========================================================================
\part{The coupled architecture}
\label{part:coupling}
% ===========================================================================

With the 3-D immobile physics working, the question became how to drive it.
The changes in this part are on the Python side and in the repository
structure; they are what turned a verification case into a simulation code.

\change{A per-node field channel replacing the four-scalar handoff}
\label{chg:fieldchannel}
The 0-D~$\rightarrow$~3-D handoff was replaced by a genuine field exchange
through MoDELib3's own cluster-dynamics configuration block.
\why{The previous route, \code{legacy/modelib\_fem.py}, wrote \emph{four
scalars} per dose step into the microstructure-generator inputs, discarding all
spatial structure. A spatially resolved code driven by a spatially averaged
boundary condition is not a coupled calculation.}
\impact{\code{modelib\_field.FieldBridge.write\_immobile\_field} writes the
per-node immobile state into the \code{evl} CD block;
\code{modelib\_qssa.MobileQSSASolver.solve} reads $C_M^*(x)$ back at the CD
nodes. \code{legacy/modelib\_fem.py} is kept, unused, only so older results
stay reproducible.}
\cost{The exchange is file-based and $O(\text{nodes})$; at 90\,617 nodes it is
a few seconds against a fast solve of several minutes.}

\change{The fast solve moved \emph{inside} the march loop}
\label{chg:fastinside}
\why{This is the load-bearing correctness point of the coupling. If the mobile
field is instead replayed from a previously recorded run, the loop is not an
operator split at all: $C_M$ never responds to the immobile state the march is
building, so a seed away from quasi-steady state can never relax. A 0-D seed
is always such a seed --- it is spatially uniform and carries no boundary
layer, while the true $C_M^*(x)$ is pinned to thermal equilibrium on every
Dirichlet face.}
\impact{\code{coupling/march.py} alternates
\code{qssa.MobileQSSASolver.solve} (DDomp with \code{useImmobileSolver=0})
with \code{run\_immobile\_step} (one OpenMP batch subprocess) every
\code{fem\_every} substeps.}
\cost{This is where essentially all the wall time goes, and it is why
Part~\ref{part:performance} exists. Splitting error is first order in the
interval between fast solves --- \code{fem\_every} substeps --- \emph{not} in
the dose-snapshot spacing, so \code{fem\_every} is the accuracy dial. It
mattered, therefore, that a value not dividing \code{substeps\_per\_interval}
was silently skipping whole intervals; the configuration validator now rejects
that.}

\change{CVODE at every quadrature point, with the mobile field frozen}
\label{chg:cvode}
The slow step integrates the immobile ODEs independently at every node with
\code{freeze\_mobile=1}, using CVODE BDF with a dense $15\times15$ linear
solver, dispatched through the 0-D solver's OpenMP \code{--batch\_file} mode.
\why{This is the design point that makes the whole scheme tractable, and it is
easy to misread the numbers as claiming something impossible. \textbf{A tightly
coupled million-equation stiff system would be hopeless. This is not one.} It
is 72\,494 independent 15-dimensional systems, not one 1\,087\,404-dimensional
one (500~nm case: 90\,617 nodes, deduplicated $\times1.25$).

The decoupling is legitimate rather than a cheat because \textbf{the only thing
coupling neighboring points is diffusion of the mobile species}, and the slow
step freezes those. Freeze the mobile field and the spatial coupling is gone
by construction: a point's immobile ODEs depend only on its own loop
populations and the mobile values pinned into its state vector.}
\impact{\code{rate\_equations.ode\_system(t, y)} takes a single point's state
vector; no neighbor state enters it. Each point is one case of the batch mode,
dynamically scheduled over 24 threads with a private CVODE workspace each.}
\cost{The cost consequence, for one implicit factorization:

\begin{center}
\begin{tabular}{lr}
\toprule
 & flops \\
\midrule
split: $72\,494 \times 15^3$ & $2.4\times10^{8}$ \\
coupled: $(1\,359\,255)^3$   & $2.5\times10^{18}$ \\
\midrule
\textbf{ratio} & $\mathbf{1.0\times10^{10}}$ \\
\bottomrule
\end{tabular}
\end{center}

\noindent
Sparsity softens the coupled figure but does not rescue it --- a 3-D FE
Jacobian still factorizes at $\sim O(N^2)$ with fill growing as $O(N^{4/3})$.

Stiffness localizes the same way. Each point's system genuinely is stiff
($\sim$100 CVODE internal steps per point per substep, which is why BDF), but
stiffness cost scales with the \emph{dimension of the coupled block}, not with
how many independent blocks there are. A $15\times15$ dense LU is
$\sim$1100~flops.

\textbf{The coupling did not vanish --- it moved to where it is affordable.}
The tightly coupled problem still exists: it is the fast solve, 362\,468
unknowns, \textbf{925~s} measured, against 28.3~s for an entire slow sweep over
90\,617 points. What the split buys is \emph{frequency}:
\code{solveMobileClusters} has no time derivative, so the coupled problem is
solved \textbf{8 times over the march} instead of at every timestep. The
counterfactual, from the same measurements ($925\,\text{s}/5$ Newton
iterations $= 185$~s per linear solve):

\begin{center}\small
\begin{tabular}{l}
$\sim$100 internal steps $\times$ 24 substeps $\div$ 5 steps per refactorization $\approx$ 480 factorizations \\
$480 \times 185\,\text{s} \approx \mathbf{25\ h}$ \quad --- the 4-species mobile block \emph{alone}, not all 19 \\
measured slow side, all 24 substeps: \quad\quad \textbf{11.3 min} \\
\end{tabular}
\end{center}

\noindent
Three things stacked: freezing the mobile field removes the spatial coupling;
the coupled part has no time derivative so it is solved 8 times rather than
thousands; and batch OpenMP plus deduplication (identical states solved once
--- $\times1.25$ here, higher early on when much of the domain shares a state)
spreads what is left over 24 cores.

The slow side is linear in node count, as this structure predicts: 9.3~s per
substep at 27\,720 nodes and 28.3~s at 90\,617 --- $3.35\times10^{-4}$ versus
$3.12\times10^{-4}$~s/node.}

\change{Configuration as a value: \code{SimulationConfig} and domain-keyed staging}
\label{chg:config}
The notebook was reduced to six control dicts, validated by
\code{dislocluster\_code.config.\allowbreak SimulationConfig.from\_dicts}.
\why{A run whose parameters live in scattered module globals cannot be
reproduced, and a validator that runs \emph{after} an expensive stage costs a
whole march.}
\impact{Unknown keys are rejected rather than ignored, and validation happens
before anything expensive runs: a dose grid that does not increase strictly, a
seed past the first snapshot, a boundary layer thicker than half the domain, an
\code{element\_order} outside $\{1,2\}$ (MoDELib's \code{SimplexReader}
consumes only msh element types 4 and 11).

Staging is keyed on the \emph{domain}, not the run: \code{cfg.domain\_key}
hashes the geometry, mesh, boundary conditions, material and temperature --- but
\emph{not} the dose grid --- so changing only the doses reuses the staged case,
its CD node set and its bootstrap.

One deliberate non-change: \code{coupling.march} still carries
\code{DOSE\_SEED}, \code{SUBSTEPS\_PER\_INTERVAL} and friends as module
globals, read by \code{default\_config()} \emph{at call time}, so that callers
which monkey-patch them still work. Turning those into dataclass field
defaults would bind them at import and make a later assignment a silent
no-op.}
\cost{Re-running a dose sweep on a staged domain skips mesh generation and
bootstrap entirely --- minutes to hours, depending on the mesh.}

\change{\code{BOUNDARY} made a real boundary-condition channel}
\label{chg:boundary}
\code{periodic\_face\_ids} goes into \code{polycrystal.txt}; the applied load
goes into \code{ElasticDeformation.txt}.
\why{Previously the only available condition was Dirichlet everywhere.}
\impact{Stress is given in \textbf{MPa} and converted on staging: MoDELib
normalizes stress by \code{mu\_SI}, so the numbers in that file are multiples
of the shear modulus, and writing MPa straight in would overstate the load by
$\sim$33\,000$\times$ for Zr.}
\cost{None.}

\change{Package consolidation and single-source path resolution}
\label{chg:paths}
All importable code was gathered into \code{dislocluster\_code/} at the
repository root, installed editable, with every repository location resolved by
\code{paths.py} from a \code{.dislocluster\_root} marker.
\why{Four modules had re-derived their locations with
\code{Path(\_\_file\_\_).parent.parent}, and that idiom broke \emph{silently}
when the package moved --- \code{cpp\_bridge} resolved its default
\code{base\_dir} that way and stopped finding \code{solver.exe} at all.}
\impact{\code{ZrMicro/py\_utils/} remains as a compatibility shim: each old
module aliases \code{sys.modules[\_\_name\_\_]} to its new home, so
\code{from py\_utils import paths}, \code{python -m py\_utils.X} and the old
notebooks all keep working, and \code{py\_utils.X is dislocluster\_code....}
is true.}
\cost{None. \code{paths.describe()} now produces an OK/MISS report per
location, and \code{paths.workbook\_drift()} hashes the two input workbook
directories and names any that have diverged --- the same failure mode that
once put the 0-D calibration five orders of magnitude off.}

% ===========================================================================
\part{Physics corrections found by measurement}
\label{part:physics}
% ===========================================================================

The changes in this part were all found the same way: a number was computed a
second, independent way, and the two did not agree.

\change{The atomic volume corrected to the hcp Zr value}
\label{chg:omega}
$\Om$ was changed from $1.2\times10^{-29}$~m$^3$ to
$\mathbf{2.326553\times10^{-29}}$~m$^3$.
\why{The old value is $V_{\rm cell}/4$ --- four atoms in an hcp primitive cell
that holds two. It implies $\sim$12.6~g/cm$^3$ for Zr, about half the true
atomic volume. Two independent derivations from data already in the input
files agree to \textbf{0.14\%}:

\begin{center}
\begin{tabular}{lr}
\toprule
route & value \\
\midrule
lattice: $(\sqrt3/2)a^2c/2$, $a = b_a = 3.23$~\AA, $c = b_c = 5.15$~\AA & $\mathbf{2.326553\times10^{-29}}$~m$^3$ \\
mass density: $M_{\rm Zr}/(\rho N_A)$, $\rho = 6520$~kg/m$^3$ & $2.3233\times10^{-29}$~m$^3$ \\
former value ($V_{\rm cell}/4$) & $1.2\times10^{-29}$~m$^3$ \\
\bottomrule
\end{tabular}
\end{center}

\noindent
The lattice route is the one adopted, so that \Om\ stays consistent with the
same $b_a$ and $b_c$ that the loop-radius prefactors $\sqrt{\Om/\pi b}$ use.}
\impact{Changed in \textbf{four} places, which must stay equal:
\code{Simulations/input/Zr\_input\_parameters.xlsx}
(\code{Physical\_Properties!D5}); \code{ZrMicro/input/} --- the same cell,
copied byte for byte so that \code{paths.workbook\_drift()} stays clean;
\code{MoDELib3/Library/Materials/Zr3d\_ghoniem.txt}
(\code{atomicVolume\_SI}); and \code{zerod/input\_data.py}, whose fallback
default read $1.4\times10^{-29}$ --- a \emph{third} number for one lattice
constant. \code{InputData.check\_atomic\_volume()} runs on every
\code{build\_sim} and is what would catch a regression.}
\cost{None computationally. But \Om\ is not a reporting factor: every number
density is $C/\Om$, line density is $2\pi r C/\Om$, and emission carries
$\exp(\sigma\Om/kT)$. Measured on the fitted set at 10~\dpa:
$N_a\ 0.483\times$, $N_c\ 0.516\times$, $c_a\ 0.522\times$, $C_v\
1.047\times$, $C_i\ 1.059\times$.}
\status{\textbf{The 28-parameter fit is now stale.} The loop densities halve,
and those parameters were fitted against experimental loop densities at the old
\Om. The model now has the right lattice constant and a fit that no longer
matches experiment. \textbf{A refit is required before comparing to data} and
is planned separately. \code{calibration.LEGACY\_OMEGA = 1.2e-29} reproduces a
pre-correction run. \textbf{Runs made before this correction are not
comparable to runs made after.}}

\change{Conservation reinterpreted as an open-system balance}
\label{chg:conservation}
The six accumulators integrate a 0-D atom balance with \textbf{no transport
term}:
$d(\text{stored})/dt = \text{production} - \text{recombination} -
\text{sink absorption}$.
On a spatially resolved march that cannot close, and the residual is
\emph{not} solver error --- it is the atoms that left through the Dirichlet
surface.
\why{Two things force it: the domains here are Dirichlet over their whole
surface unless faces are made periodic, and the slow step \emph{freezes} the
mobile species while production keeps accumulating into them (the mobile pool's
own balance is closed by the \emph{fast} solve, which is where the flux to the
boundary lives).

The magnitude settles it: the 200~nm march reports a residual of \textbf{$-71$\%
of cumulative production}. Nothing in a CVODE run at \code{rtol=1e-6} is 71\%
wrong.}
\impact{\code{post\_process.\_calculate\_conservation} also publishes
\code{grain\_boundary} ($= -$residual, positive when atoms have left), plus
\code{mobile\_change} and \code{immobile\_change}.
\code{post.volume\_average} presents it as a physical channel because it
\emph{knows} the surface is absorbing --- it reads \code{periodic\_face\_ids}
from \code{config.json}. A standalone 0-D run leaves it as the numerical
residual it genuinely is there (closing to $3\times10^{-12}$), and
\code{plot\_conservation} falls back to the old relative-error figure.}
\cost{None. It is closure \emph{by difference}, not an independent
measurement --- it inherits every other channel's error, and its credibility
rests on those being exact solver integrals rather than reconstructions.
Which is why Change~\ref{chg:boundaryflux} exists.}

\change{Conservation figures in defect counts, not fractions}
\label{chg:counts}
The fate figures were changed from fractions of production to absolute defect
counts, with $n_{\rm atoms} = V/\Om$, $V$ from \code{fields.domain\_volume} ---
the crystal, not its bounding box.
\why{A fraction hides the size of the system it is a fraction of, and the
\code{*\_fractions} figures had a ``Sum (=1)'' line that was reaching
$\sim$0.3, because the boundary channel was missing from it.}
\impact{The 200~nm prism is $6.93\times10^{8}$ atoms, and at 10~\dpa:

\begin{center}
\begin{tabular}{lrr}
\toprule
channel & interstitials & vacancies \\
\midrule
produced & $6.24\times10^{9}$ & $6.24\times10^{9}$ \\
recombined & $1.43\times10^{9}$ (22.9\%) & $1.43\times10^{9}$ (22.9\%) \\
absorbed at network sinks & $3.51\times10^{8}$ (5.6\%) & $3.20\times10^{8}$ (5.1\%) \\
\textbf{absorbed at grain boundary} & $\mathbf{4.45\times10^{9}}$ (71.4\%) & $\mathbf{4.48\times10^{9}}$ (71.9\%) \\
stored in microstructure & $4.71\times10^{6}$ (0.1\%) & $1.33\times10^{6}$ (0.0\%) \\
\bottomrule
\end{tabular}
\end{center}

\noindent
Equal Frenkel production and equal recombination are exact; the storage
asymmetry ($4.7\times10^{6}$ against $1.3\times10^{6}$) is the bias-driven
excess that drives growth. The ``Sum (=1)'' line is true again.}
\cost{None.}

\change{The accumulators are per-interval, not cumulative}
\label{chg:peraccum}
\code{volume\_average.cumulate\_accumulators} was added to cumulative-sum the
six accumulator columns along the dose axis before they are reported.
\why{The march restarts each interval's integration from the marched state, so
each snapshot's accumulators hold that \emph{interval's} contribution, not the
running total. Reporting them as if cumulative understates every channel by the
history preceding the last interval.}
\impact{\code{build\_results(..., per\_interval\_accumulators=True)}.}
\cost{None.}
\status{This was found \emph{by} the independent boundary-flux check of
Change~\ref{chg:boundaryflux} --- exactly the job that check was added to do.
An earlier hypothesis, that the discrepancy came from mixing the old and new
\Om, was tested with \code{LEGACY\_OMEGA} and \textbf{was wrong}: it barely
moved the ratios.}

\change{The grain-boundary channel measured a second, independent way}
\label{chg:boundaryflux}
\code{post/boundary\_flux.py} obtains the boundary flux without ever touching
the accumulators. The fast solve returns a \emph{steady} mobile field, so for
each mobile species $0 = \nabla\cdot(D\nabla C) + R$, and by the divergence
theorem
\[
  \int_\Omega R \,dV \;=\; \oint_{\partial\Omega} (-D\nabla C)\cdot n \,dA
  \;=\; \text{net outward flux}.
\]
\why{Change~\ref{chg:conservation} is closure by difference. A channel obtained
by subtracting everything else from production is only as good as everything
else, and a 71\% number deserves a second opinion.}
\impact{$R$ is exactly the 0-D right-hand side for the mobile components ---
that model has no diffusion, so its mobile equations \emph{are} the reaction
terms, and the grain boundary appears nowhere in it. That is why the volume
integral of its mobile RHS measures the flux to the boundary. Atom weights are
$(1,2,3)$ for $(C_i, C_{2i}, C_{3i})$ and 1 for $C_v$, matching the stored
channels.

Measured against closure on the 200~nm march:

\begin{center}
\begin{tabular}{rrrr}
\toprule
dose & measured & closure & ratio \\
\midrule
0.01 & 5\,666\,311 & 5\,106\,945 & 1.110 \\
0.1  & 44\,999\,294 & 44\,632\,725 & 1.008 \\
1    & 431\,455\,063 & 444\,694\,897 & 0.970 \\
10   & 4\,287\,277\,443 & 4\,450\,899\,303 & \textbf{0.963} \\
\bottomrule
\end{tabular}
\end{center}

\noindent
\textbf{Agreement to 3.7\% at 10~\dpa} --- an independent confirmation, not an
identity. The 11\% at 0.01~\dpa\ is the trapezoid over a coarse log grid where
the rate falls 36\% across the first interval, not a physics discrepancy.}
\cost{Seconds. \code{march\_report} calls it \emph{guarded}, so a diagnostic
can never cost a report that a multi-hour march earned.}
\status{A surface integral inside MoDELib, where the FE gradients already
exist, would be the gold standard. This needs no C++ change and works on runs
that already exist, which is why it came first. On the 500~nm reference march
the ratio sits at a stable 0.63--0.65; the leading hypothesis is
Part~\ref{part:open}, item~4.}

% ===========================================================================
\part{Computational efficiency}
\label{part:performance}
% ===========================================================================

Everything in this part leaves the answer unchanged. Where that was checkable
bit for bit, it was checked bit for bit, and the check is reported alongside
the speedup --- a speedup obtained by changing the answer is not a speedup.

% ---------------------------------------------------------------------------
\section{The 0-D solver}
% ---------------------------------------------------------------------------

\change{An exact Jacobian by forward-mode automatic differentiation}
\label{chg:ad}
The right-hand side was moved into a single scalar-type-templated core, so that
the residual and an \emph{exact} Jacobian are two instantiations of one
definition. SUNDIALS' difference-quotient Jacobian was retired.
\why{A difference-quotient Jacobian on a system spanning nine orders of
magnitude in concentration is inaccurate where it matters most, and it costs
$n$ extra RHS evaluations per Newton matrix.}
\impact{One definition, two instantiations --- so the Jacobian cannot drift out
of step with the residual, which is the usual failure mode of a hand-written
one.}
\cost{At $n = 19$ the dense linear algebra was never the bottleneck, so this
bought accuracy and robustness rather than wall time. Reported as measured.}

\change{The six accumulators moved out of the implicit system}
\label{chg:quadrature}
The conservation accumulators were recognized as pure quadratures and moved
into CVODES quadrature variables.
\why{They are integrals of quantities the solver already computes; nothing
depends on them, so they need not be inside the Newton block.}
\impact{The Newton block shrinks from $19\times19$ to $13\times13$ standalone.
Under the operator split the integrated width is \textbf{15}, not 9: with the
mobile species frozen, \code{acc\_mode=2} (\code{ACC\_STATE\_RELAX},
\code{N\_RLX\_FROZEN=15}) is what \code{red\_dim()} returns. The narrower
$9\times9$ form is \code{acc\_mode=1}, which is not used --- it measured
\emph{slower}.}
\cost{Modest at this size, and honestly so: at $n = 19$ dense linear algebra is
not where the time goes.}

\change{One SUNDIALS workspace per OpenMP thread}
\label{chg:workspace}
Solver objects are constructed once per thread and reused across cases instead
of being constructed and destroyed per case.
\why{Profiling with the integration window collapsed to zero showed
per-case construction and destruction of SUNDIALS objects consuming
$1.31$~s of thread time against $0.2$~s of actual solving in a 256-point
coupling substep. The measurement redirected the entire effort.}
\impact{A thread-local workspace pool in the batch driver.}
\cost{\textbf{$1.67\times$ on end-to-end march wall time} at $q = 1024$
quadrature points --- more than the two changes above combined. This is the
clearest case in the project of measurement beating expectation: the two
changes that were requested for speed delivered accuracy, and the one that
delivered speed was not on the list until the profile was read.}

% ---------------------------------------------------------------------------
\section{The fast solve}
% ---------------------------------------------------------------------------

The fast solve is the tightly coupled part of the split and therefore where
essentially all the wall time is: on the 500~nm reference march, 1~h~05~m of
fast solves against 12~m of slow substeps. It was instrumented before it was
touched, and \code{compute()} --- the preconditioner factorization --- was
found to hold 85--87\%.

\change{Skip the elastic solve when it can only return zero}
\label{chg:elastic}
The fast step now sets \code{useElasticDeformation=0} and
\code{useElasticDeformationFEM=0} when the staged case has no applied load and
no dislocations.
\why{\code{useElasticDeformation=1} costs $\sim$42~s per step on a 72\,345-DOF
solve for a case with no dislocations and no applied stress --- roughly 20
minutes of a 30-step run, spent computing a field that is identically zero.}
\impact{\code{staging.inputs.\allowbreak elastic\_is\_trivial(sim\_dir)} returns
\code{(trivial, reason)}, and \code{qssa.MobileQSSASolver} takes an
\code{elastic=} argument. The \emph{reason} is returned rather than a bare
boolean so that a run which unexpectedly keeps the elastic solve says why.}
\cost{\textbf{$1.86\times$} on the 500~nm case, bit-identical.}

\change{Remove \code{AcIR}, a matrix assembled every iteration and never read}
\label{chg:acir}
\why{Dead work in the innermost loop.}
\impact{One deletion in \code{ClusterDynamicsFEM.cpp}.}
\cost{\textbf{$1.64\times$} --- the single largest of the four fast-solve
changes, and it removed code rather than adding any.}

\change{Cache the Dirichlet selection matrix; build $A_1$ by filtering}
\label{chg:filter}
The selection matrix $T$ and its index map are cached, and $A_1 = T^{T}AT$ is
replaced by a filtered copy of $A$'s nonzeros.
\why{$T$ is a \emph{selection} matrix, so $T^{T}AT$ is a submatrix --- forming
it as two sparse products is doing general-purpose work to accomplish a
filter.}
\impact{\code{rMap} added to \code{FixedDirichletSolver.h}, along with a new
\code{computeFromTriplets\allowbreak <BilinearWeakFormType>()}.}
\cost{\textbf{$\sim$1.00$\times$ alone} --- it measured as nothing, because
\code{rSolver} was rebuilt every iteration and discarded the cache before it
could pay. Reported at its measured value; Change~\ref{chg:hoist} is what made
it count.}

\change{Parallelize the element assembly}
\label{chg:parallel}
\code{BilinearWeakForm::globalTriplets} parallelized over explicit contiguous
chunks, concatenated in chunk order.
\why{Element assembly is embarrassingly parallel and was serial.}
\impact{Chunks are explicit and contiguous, and concatenated \emph{in chunk
order}, specifically to preserve the serial triplet sequence.
\code{setFromTriplets} sums duplicate entries, so a different triplet order
re-associates those sums and the result would agree only to rounding --- which
would have destroyed bit-identity as a check. Keeping the order is what makes
the verification possible.}
\cost{\textbf{$1.33\times$}.}

\change{Hoist the loop-invariant weak forms out of the Newton iteration}
\label{chg:hoist}
\code{R1c}, \code{bWF\_R1c}, \code{RIc}, \code{bWF\_RIc}, the triplet head and
tail, and \code{rSolver} were moved above the Newton loop; each iteration
splices head~+~$R_2$~+~tail.
\why{Only the $R_2$ block depends on the current iterate. Everything else was
being rebuilt from scratch every iteration.}
\impact{Per-iteration work reduces to assembling one block and splicing it
between two cached triplet ranges.}
\cost{\textbf{$1.11\times$} on its own --- and it is what finally paid off
Change~\ref{chg:filter}, since \code{rSolver} now survives across iterations.

\medskip\noindent
\textbf{Cumulative, Changes~\ref{chg:acir}--\ref{chg:hoist}:} on one isolated
DDomp call at 200~nm, $96.6\,\text{s} \rightarrow 37.8\,\text{s}$,
\textbf{2.556}$\times$, with the CD block $(27\,720, 12)$
\textbf{bit-identical}.}

\change{The optimized fast solve measured on a full march}
\label{chg:abcompare}
Two complete 500~nm hexagonal marches were run back to back on an otherwise
idle machine, differing \emph{only} in the DDomp binary.
\why{An isolated-call speedup is not a march speedup: Amdahl applies, and the
slow side is untouched.}
\impact{\code{studies/compare\_architectures.py} reports the physics first and
the timings second, and uses \textbf{the slow step as a control} --- it runs the
same code in both runs, so a ratio far from 1.0 there means the two runs were
not measured under comparable load and the fast-side number must not be
quoted.}
\cost{
\begin{center}
\begin{tabular}{lrrr}
\toprule
quantity & reference & optimized & speedup \\
\midrule
march total & 1~h 18~m 05~s & 42~m 41~s & \textbf{1.829}$\times$ \\
fast solves ($8\times$ DDomp) & 1~h 04~m 53~s & 29~m 49~s & \textbf{2.176}$\times$ \\
one fast solve & 8~m 06~s & 3~m 43~s & 2.176$\times$ \\
slow substeps (control) & 12~m 18~s & 11~m 58~s & 1.028$\times$ \\
\bottomrule
\end{tabular}
\end{center}

\noindent
The state array $(9, 90\,617, 19)$ is \textbf{bit-identical} between the two
runs, and the control ratio of 1.028 confirms the machine was comparably
loaded. The march speedup is below the isolated 2.556$\times$ for two reasons,
both expected: the slow side is untouched, and the first interval takes no fast
solve at all.}

% ===========================================================================
\part{Portability, reproducibility and post-processing}
\label{part:support}
% ===========================================================================

\change{One build script, one \code{CMakeLists}, every platform}
\label{chg:build}
\why{The tree only built on the platform it was last built on. The build
script \code{sed}-patched \code{MoDELib3/CMakeLists.txt} on every run ---
rewriting the Eigen path, deleting the Qt lines, commenting out
\code{add\_subdirectory(DDqt)} --- so the source tree carried whichever
platform's edits had been applied last. On macOS the build stopped at the first
translation unit.}
\impact{Four platform-specific failures, none of which Linux ever shows:

\begin{center}
\small
\begin{tabular}{p{6.2cm}p{8.4cm}}
\toprule
Symptom & Cause \\
\midrule
\code{clang: error: unsupported option '-fopenmp'} & the flag was hard-coded;
  Apple clang needs \code{-Xpreprocessor -fopenmp} and Homebrew's libomp, which
  \code{find\_package(OpenMP)} now supplies \\
\code{use of undeclared identifier 'assert'}, 103 files & those files never
  included \code{<cassert>}; libstdc++ leaks it in, libc++ does not. A
  tree-wide \code{-include cassert} substitutes for editing all 103 and keeps
  the fork mergeable \\
\code{must explicitly initialize the const member 'invTrD'} & Eigen
  $\ge$ 3.4.90 declares \code{Matrix() = default}, so a const member with no
  initializer is ill-formed. \code{invTrD} is unused; now zero-initialized.
  An Eigen-version dependency, not a platform one \\
\code{no viable conversion from 'RationalLatticeDirection<3>'} & clang resolves
  the dependent \code{this->operator+(...)} call against the
  \code{LatticeVector} overload alone; naming the temporary fixes it \\
\bottomrule
\end{tabular}
\end{center}

\noindent
Every dependency is now searched for, every optimization flag probed with
\code{check\_cxx\_compiler\_flag} before use, and everything optional degrades
to a status line. \code{DDqt} is built when Qt6 and VTK $\ge$ 9.4 are both
found and skipped otherwise. \code{cmake\_minimum\_required} went from 3.1.0 to
3.16 --- CMake $\ge$ 4.0 refuses to run a project asking for less than 3.5, so
the old floor was a time bomb.}
\cost{None. \code{DDomp}, \code{microstructureGenerator} and
\code{libMoDELib.a} all build clean on Apple silicon, and the CD Newton
iterates reproduce \emph{bit-identically} between two differently configured
builds of the same tree.}

\change{Eigen 3.4.x required, and the version test made fatal}
\label{chg:eigen}
\why{This is the worst kind of portability failure: a clean build and a wrong
answer, reported as a solver problem. Homebrew's \code{eigen} formula is 5.x,
and MoDELib built against it \emph{compiles} --- then the BiCGSTAB solve inside
the mobile Newton iteration breaks down on the first step
(\code{Iterative FixedDirichletSolver failed}) on a 500~nm case that had staged
successfully on Linux/Eigen 3.4.0 two months earlier. Rebuilding the same
source against 3.4.0 and re-running the same case bootstraps in 138~s.}
\impact{\code{build\_modelib.sh} fetches 3.4.0 into
\code{Libraries/eigen-3.4.0} when the system has only Eigen 5;
\code{CMakeLists.txt} prefers that copy; the version check is \emph{fatal},
with \code{-DEIGEN3\_ALLOW\_UNTESTED=ON} to override.

The version test has to read \emph{both} places Eigen states its version:
3.4.x puts the macros in \code{Eigen/src/Core/util/Macros.h}, 5.x in
\code{Eigen/Version}, and 5.0.1 spells itself \code{3.5.0} there
(\code{EIGEN\_WORLD\_VERSION} stays 3 forever) --- so a naive
\code{\textasciicircum3\textbackslash.} test passes it on both counts.}
\cost{None.}

\change{No build tree in git; \code{preflight} runs the binaries}
\label{chg:preflight}
\why{\code{MoDELib3/build\_dc/} --- \code{libMoDELib.so}, \code{DDomp}, and a
\code{build.ninja} carrying 374 absolute \code{/mnt/d/...} paths --- used to be
committed. A clone on another machine got that Linux ELF,
\code{paths.modelib\_ddomp()} reported MoDELib as built, \code{ensure\_modelib}
skipped the build, and the first DDomp call failed.}
\impact{Both build trees are ignored, and \code{build.preflight()} \emph{runs}
each binary rather than trusting its presence.

A related silent failure: moving or renaming the repository invalidates both
build trees, because a CMake cache stores absolute paths. The failure is quiet
--- \textbf{the already-built binary keeps working}, so nothing breaks until the
next reconfigure, possibly months later. \code{build.stale\_cache\_home()}
reports it, \code{build.describe()} flags it, and
\code{ensure\_zrmicro\_solver()} discards and rebuilds. That check runs
\emph{before} the ``already built'' shortcut --- putting it after, as an earlier
revision did, made it unreachable in exactly the situation it exists for.}
\cost{Seconds per run, against hours of a wrong one.}

\change{\code{evl/cdNodes.txt} --- the CD node coordinates}
\label{chg:cdnodes}
\why{Without it the fields cannot be placed in space: the CD trial functions
live on second-order elements (24\,115 nodes for a 3392-node mesh), so the
field rows do \emph{not} correspond to the \code{.msh} vertices and the
ordering is not reproducible externally.}
\impact{Written once per run.}
\cost{Negligible.}

\change{Geometry-true field rendering}
\label{chg:figures}
\why{Nothing in the figure code should be specific to one domain shape, and
several things were.}
\impact{Every shape decision --- the outline, the hidden-line test, the volume
that turns a density into a count, the region loops are placed in, the clip ---
comes from \code{fields.domain\_faces} / \code{domain\_edges} /
\code{domain\_volume}, the convex hull of the CD node cloud. A cubic case gives
6 faces, 12 edges and \code{domain\_volume == prod(hi - lo)} to machine
precision, so the box case is untouched; hexagonal gives 8 faces and 18 edges.

Three specific corrections, each of which had been silently wrong on a
non-cubic domain:

\begin{itemize}\itemsep2pt
\item \emph{Both cut planes go into ONE \code{Poly3DCollection}.} Matplotlib
  depth-sorts 3-D \emph{artists} by a single scalar each, so two intersecting
  surfaces drawn as separate \code{plot\_surface} calls cannot interleave and
  whichever sorts in front hides the other \emph{completely} --- which is what
  used to happen, the vertical cut covering the horizontal one entirely.
\item \emph{The panel zoom is fitted, not fixed.} A 3-D axes clips its artists
  at the axes rectangle, not at the data cube, so a zoom that projects the
  domain taller than the panel silently cuts the ends off the wireframe. It is
  the \emph{aspect ratio} that decides this, not the size: a cube fits at
  $1.22$ and is left exactly as it was, while a $1:0.87:1.6$ hexagonal prism
  projects to $-0.068..1.054$ and loses both end faces --- identically at
  200~nm and at 500~nm. \code{\_fit\_zoom} reduces the zoom until the convex
  \emph{hull} fits, measured through matplotlib's own \code{get\_proj}. Fit to
  the hull, not the bounding box, whose corners for a prism sit out in the six
  empty wedges.
\item \emph{The volume that turns a density into a count} is
  \code{domain\_volume}, not $\prod(\text{bounding box})$ --- for a hexagonal
  prism the box is $4/3$ the crystal, so the discrete loop count used to come
  out a third too high.
\end{itemize}}
\cost{The loop-overlay panels are the slow figures at $\sim$16~s each; the rest
are seconds.}
\status{\code{volume\_average.averaged\_trajectory} still calls
\code{voronoi\_weights} \emph{without} \code{faces}, so on a non-box domain its
volume weights carry the same bounding-box error. That is deliberate --- fixing
it changes already-published figures --- but it is wrong for a hexagonal run.
See Part~\ref{part:open}, item~4.}

\change{Interpolated movie frames}
\label{chg:movies}
\code{OUTPUT['movie\_interp']} (default 5, \code{1} disables) subdivides each
dose interval into that many frames.
\why{The march writes \emph{one} snapshot per dose interval, so an eight-dose
run animates as an eight-frame flipbook however high the fps.}
\impact{\textbf{The extra frames are interpolated, not solved}, and every
synthetic frame is labelled \code{(interp)} in its title. Two things about the
blend are load-bearing:

\begin{itemize}\itemsep2pt
\item It is \emph{geometric} wherever both endpoints are positive (linear
  otherwise --- the pristine seed sits at dose 0). The fields grow by decades
  through the nucleation transient, where a linear interpolant would sit at the
  upper endpoint for the whole interval, and geometric blending commutes with
  the power-law reductions the figures draw: a size $d \sim (c/N)^{1/3}$ formed
  from blended $c$ and $N$ is the blended $d$ to $2\times10^{-15}$.
\item It is taken on the \emph{CD block}, not on the 19-column state it is
  built from. \code{immobile\_0d\_to\_modelib} sums species, and a sum of
  geometric blends is not the geometric blend of the sums, so blending the
  state first yields frames that are \emph{not bracketed by their neighbors}
  --- visible as a population overshooting and falling back.
\end{itemize}

\noindent
\code{cd\_blocks} deliberately keeps \code{interp=1} as its default:
\code{discrete\_loops} and the \code{3d/}/\code{gb/} panels must see solved
states only. Only \code{movies.render} opts in.}
\cost{Interpolation is free relative to solving; it is what makes a smooth
movie possible without more fast solves.}

\change{Computational statistics in the run provenance}
\label{chg:stats}
\code{march\_report.computational\_stats} writes wall clock total and per
solve, the number of ODEs actually integrated, node counts and per-node rates
into \code{report.md}.
\why{Cost claims about this architecture were being made from memory. They
should come from the run.}
\impact{Reports the \emph{integrated} ODE count, not the carried state length
--- 15 per node under the operator split, not 19.}
\cost{Negligible.}

% ===========================================================================
\part{Anisotropic diffusion and the discrete handoff}
\label{part:aniso}
% ===========================================================================

Work of 16 August 2026. Two threads: making the diffusion tensor and the
capture efficiencies describe the same physics, and building the
continuum\,$\rightarrow$\,discrete transition listed as open item~1 below.

\change{One anisotropy factor per species, generating both representations}
\label{chg:aniso-one-source}
The anisotropy appeared in the material file twice --- as
\code{mobileSpeciesEnergyMigration\_eV}, the tensor the fast solve diffuses
with, and as \code{dadAnisotropy}, the $p_m$ the closed-form capture
efficiencies use. They were independent. \code{staging/anisotropy.py} now
generates both from one $p_m$ per species, split at fixed
$D_{\text{eff}}=(D_a^2D_c)^{1/3}$:
\[
E_m^{\langle 11,22\rangle}=E_m^{\text{eff}}+2k_BT\ln p_m,\qquad
E_m^{\langle 33\rangle}=E_m^{\text{eff}}-4k_BT\ln p_m .
\]
\why{The code simultaneously believed diffusion was isotropic and that the
sinks were biased by its anisotropy. Those cannot both be true.}
\impact{Reproduces Table~2 of Li et al.\ to the last quoted digit. Because the
anisotropy is stored as \emph{energies}, $p_m(T)=\exp(-\Delta E_m/6k_BT)$ is
temperature dependent, repairing a fit previously frozen at one temperature.}
\cost{Turning the tensor on costs about $2\times$ in the fast solve (223.7\,s
isotropic against 433.8 and 466.1\,s anisotropic, fresh runs on one mesh). The
\emph{degree} of anisotropy does not matter: $p=0.70$ and $p=0.9137$ cost the
same despite a factor~5 difference in directional contrast.}

\change{A closed-form criterion for simultaneous growth of both loop families}
\label{chg:cogrowth}
Both growth conditions reduce to bounds on one number, the arrival ratio
$A/B=\bar D_v c_v\big/\sum_m \bar D_m c_m|m|$:
\[
\frac{Z^0_I}{Z^0_v}\frac{\bar p}{p_v} \;<\; \frac{A}{B} \;<\;
\frac{Z^0_I}{Z^0_v}\frac{\bar f}{f(p_v)},\qquad f(p)=\tfrac12(p+p^{-2}),
\]
and the window is non-empty iff $g(p_I)<g(p_v)$ with $g(p)=2/(1+p^{-3})$
strictly increasing --- that is, \textbf{iff $p_I<p_v$}.
\why{The observation the DAD mechanism exists to explain can be derived rather
than searched for. The sink strength $S_k$ multiplies both channels and cancels
from the \emph{sign}, so the microstructure sets the rate and never the
direction.}
\impact{Confirmed 16/16 on a two-grid sweep. The measured co-growth set is
$p_m=(1,\,0.9137,\,0.9137,\,0.9137)$ --- vacancies isotropic, interstitials at
the already-fitted value --- giving both families positive absorption on 100\%
of interior nodes at every dose from 0.01 to 10\,dpa, tolerance
$p_v\in[0.94,1.07]$. Li et al.'s structure gives none here: the clusters carry
3.1\% of the interstitial arrival and an isotropic species is inert in the
criterion.}
\status{The criterion governs the \emph{absorption flux} only. The content
update also carries cascade nucleation $G_k$, which no anisotropy affects, so
it does not govern net population evolution. An isotropic march is not a zero
of the criterion but another parameter point, so a ratio against one cannot
test a statement about signs.}

\change{The $d_{\text{coarsen}}$ detector, resolved to a substep}
\label{chg:dcoarsen}
\code{post/coarsening.py} forms the Avrami overlap fractions per interior node
per family and reports the dose at which they cross $\varphi^\ast$;
\code{coupling/march.py} feeds it inside the immobile substep loop and records
the trajectory in \code{summary.json}.
\why{Interpolating the crossing across a decade of dose bracketed it only to a
factor of three.}
\impact{Three independent 500\,nm marches report $d_{\text{coarsen}}=0.4$\,dpa
for the basal family and no crossing for any prismatic family, identically.}

\change{The runtime continuum\,$\rightarrow$\,discrete transition and its ledger}
\label{chg:transition}
\code{coupling/transition.py} converts a family's continuum field into discrete
loops from a live march state, hands them to DD, and removes from the continuum
only what DD can take. Armed by \code{COUPLING['discrete\_transition']}, off by
default.
\why{Once a loop is discrete it is a sink through the Green's function; if it
is also still in the continuum it absorbs twice. The same bug has a second
face --- \code{DislocationQuadraturePoint::cCD} samples the continuum field at
the loop line, so such a loop sees its own depletion twice over.}
\impact{Three leaks were found by the ledger rather than by inspection.
(i)~\code{microstructureGenerator} silently refuses loops whose nodes leave the
grain; \code{fits\_in\_crystal} now predicts that decision exactly and keeps
the refused share in the continuum. The share must be counted in
\emph{defects}: one refusal of six was 83\% by count and 98\% by defects.
(ii)~\code{sample\_family} draws an integer count from a real expectation, so
the drawn population missed the continuum total by 15.4\% at six loops;
\code{rescale\_to\_defects} fixes it exactly. (iii)~Sink strength jumps at the
handoff by $\sqrt{b_{cd}/b_{dd}}=1.414$ times $1/\text{loopSinkScale}=3.43$,
and coalescence moves it a further $5.3\times$. Conservation of stored defects
now holds to $10^{-16}$.}
\status{The sink jump is a modelling decision, reported and never asserted on.
It must be settled before a transfer feeds a solve. A transfer also requires a
domain that can hold the loops: the kept fraction is 0.997 at 0.1\,dpa on the
500\,nm case, 0.55 from 1\,dpa on, and exactly 0 on the 200\,nm case, where the
basal radii reach 65\,nm against a 173\,nm extent.}

\change{A screened cutoff on the Galerkin climb assembly}
\label{chg:climb-cutoff}
New optional material key \code{climbNeighborCutoff\_b} truncates
\code{GalerkinClimbSolver}'s $O(N_{\text{seg}}^2)$ pair assembly, which is
100\% of the climb cost --- the live path is lumped and its ``solve'' is one
scalar division per node.
\why{The kernel is not bare $1/r$ but screened as $e^{-kr}/r$ with the same
$k^2$ \code{ImmobileSinks.h} assembles; and the continuum already carries the
mean-field response, so the discrete sum must supply only the near-field
correction the mean field misses.}
\impact{Zero or absent means no cutoff, so every existing material file is
unaffected. Verified bit-for-bit: an enormous cutoff reproduces the uncut
assembly exactly in \code{evl} and in every \code{F} field. The measured
truncation error on $\dot\beta^P_{33}$ is 19.3\,/\,9.6\,/\,3.9\,/\,0.25\% at
$R_c/L_s=1/2/3/4$, so the default is \textbf{4, not the 3 originally
proposed}.}
\cost{$274\times$ fewer pairs than the naive assembly on the measured case, of
which $28\times$ is the side count and $9.6\times$ the cutoff itself.}

\change{Render side count separated from the export side count}
\label{chg:dd-sides}
\code{sides} is now the rendering polygon count and \code{dd\_sides} the one
\code{write\_microstructure} exports, fixed at 12 for the basal family.
\why{Raising the basal family to 64 sides to draw it as a circle also changed
what was exported, and the climb assembly is $O(N_{\text{seg}}^2)$.}
\impact{$28.4\times$ in the solve, for a shape difference no solve can resolve.
\code{polygon\_circumradius} matches area at either count, so the stored defect
number is identical to every digit and no published figure moves.}

\change{The prismatic-loop ellipse reduced-order model}
\label{chg:ellipse-rom}
\code{coupling/ellipse\_rom.py}: two ODEs per loop, the curvature self-stress
that gives the ROM a steady shape at all, and the capture-diffusivity closure
as an explicit choice.
\why{A prismatic loop that stays convex and centred is two numbers rather than
$2N_{\text{sides}}$, and it is the family where that holds ($2r/d=0.091$).}
\impact{Validated against what it must satisfy by construction: the defect
count reproduces \code{discrete\_loops}' radius rule exactly, the steady aspect
ratios are 2.915 and 8.500 for the two closures at $p=0.7$, and isotropic $D$
gives a rate ratio of exactly~1.}
\status{Two qualifications found by integrating it. The \emph{sign of growth}
decides which axis is major, so conclusion~2 of Li et al.\ is a statement about
growing loops; at the arrival ratio this model carries, the prismatic family is
shrinking and the ROM predicts the opposite ellipticity. And the attractor is
approached slowly --- a loop must grow some $10\times$ its nucleus to come
within 7\% of $p^{-3}$ --- so observed ellipticity measures growth history as
much as it measures $p$.}

\change{\code{COUPLING} keys that validated and were then discarded}
\label{chg:config-forward}
\code{SimulationConfig.from\_dicts} builds \code{Coupling} from an explicit
field list, and \code{phi\_star}, \code{frac\_star} and \code{coarsen\_hold}
were not in it.
\why{They were declared on the dataclass and accepted by \code{\_take}, then
dropped, so \code{COUPLING['phi\_star']} silently did nothing and every run used
the dataclass default.}
\impact{A key that validates and is then ignored is worse than an unknown one,
which \code{\_take} rejects loudly. All six coupling keys now forward. No
result is invalidated --- every run used $\varphi^\ast=0.15$, which is what was
intended --- but the parameter was only ever actually varied through the
sensitivity sweep in \code{post/coarsening.py}.}

\change{Non-ASCII status glyphs no longer kill a run on Windows}
\label{chg:stdio}
49 \code{print} lines across 12 modules carry \code{U+2713} and similar. On
Windows both a console and a \emph{redirected} stream get cp1252, which encodes
none of them.
\why{A march died at \code{input\_data.py:38} after 914\,s of CD node
bootstrapping and before a single dose step. Earlier marches survived only
because a \emph{resumed} march skips \code{driver.prepare}, which is the call
that reaches \code{InputData}.}
\impact{Fixed at the stream in \code{dislocluster\_code/\_\_init\_\_.py}, once
at import, because fixing the 49 call sites would not hold --- the next status
marker anyone types reintroduces it. A no-op on Linux and macOS.}

% ===========================================================================
\part{Open items}
\label{part:open}
% ===========================================================================

These are unresolved as of 16 August 2026. Two of them are physics decisions
rather than defects.

\begin{enumerate}\itemsep4pt

\item \textbf{The continuum\,$\rightarrow$\,discrete loop transition ---
\emph{built}, not yet exercised over a full march.} This was the largest
capability the two branches did not share. It now exists, by a different route
from fullCD's \code{initializeDiscreteClimbLoops()}: the detector of
change~\ref{chg:dcoarsen} decides \emph{when}, \code{coupling/transition.py}
performs the handoff with the conservation ledger of
change~\ref{chg:transition}, and the cutoff of change~\ref{chg:climb-cutoff}
makes the climb solve affordable. What remains is to run it: the pieces are
validated individually --- conservation to $10^{-16}$, the cutoff bit-for-bit
against the uncut assembly --- but a march that fires the transition at
$d_{\text{coarsen}}$ and accumulates growth strain through \code{F} across the
whole dose range has not yet been completed and checked. Until it has, the
sink-strength discontinuity at the handoff remains an open modelling decision
rather than a settled one.

\item \textbf{The mobile Newton iteration does not converge from a cold start.}
\code{convergenceError} $= 2.9\times10^{5}$ after the first iteration, then the
BiCGSTAB solve fails on the second. The cause is the cold start:
\code{mobileClusters} is initialized to \emph{thermal equilibrium}
($\sim10^{-15}$) and must reach the irradiation steady state ($\sim10^{-6}$) ---
nine orders --- in one Newton solve of a quadratically nonlinear system.
Switching \code{rSolver} to the direct SparseLU branch was tried and rejected:
it fails \emph{earlier}, in \code{compute()} itself, on the $\sim$96k-DOF
matrix. Note that $dt$ does not enter the mobile equation at all, so ramping
the dose step cannot help; the remedy has to be a better initial guess --- which
is exactly what the two-time-scale coupling provides --- or continuation in
$G$. Change~\ref{chg:omegacancel} removed a factor of 3.98 on the diffusion
operator after this was diagnosed, so it should be re-tested.

\item \textbf{$C_{2i}$/$C_{3i}$ depleted $\sim$2.5$\times$ once the loops have
grown.} At 1~\dpa\ they agree with the 0-D to 1--2\%; by 5~\dpa\ they have
fallen to $0.41$--$0.45\times$ and stay there. ZrMicro's loops absorb $2i$ and
$3i$ for growth --- $\text{flux}_i = \omega_i(C_i + 2C_{2i} + 3C_{3i})$ --- but
its $dC_{2i}/dt$ and $dC_{3i}/dt$ never debit those pools. Zr3d\_ghoniem's
\code{ImmobileSinks} \emph{does} debit them, so the 3-D carries a $2i$/$3i$ loop
sink the 0-D lacks. The one-line estimate
$C_{2i}^{3D}/C_{2i}^{0D} = \Lambda_{0D}/(\Lambda_{0D}+D_ik^2_{\rm loops})$
gives $20.92/(20.92+17.86) = 0.54$ against the $0.41$ measured --- same
mechanism, right order.

\textbf{Not fixed, because it is a physics decision, not a bug in either code.}
Zr3d\_ghoniem is the mass-conserving one: a loop that absorbs a
di-interstitial must remove it from the mobile pool. The reconciling change
belongs in ZrMicro, but it would shift the 0-D fit that the experiments are
calibrated against.

This matters beyond itself: the clustering nucleation flux of
Change~\ref{chg:clustering} is $K C_i C_{3i}$, so a depleted $C_{3i}$ feeds it
directly. The 3-D's own nucleation rate is $0.782\times$ the 0-D's for this
reason, and $0.704/0.782 = 0.90$ --- i.e.\ with the cluster concentrations
matched, the \aloop\ density would land within 10\%. \textbf{Closing this
should close most of what remains on the \aloop\ families as well.}

\item \textbf{The boundary-flux ratio sits at 0.63--0.65 on the 500~nm march},
against 0.963 at 10~\dpa\ on the 200~nm case (Change~\ref{chg:boundaryflux}).
The leading hypothesis, untested, is the bounding-box Voronoi weights in
\code{averaged\_trajectory} noted under Change~\ref{chg:figures}: they
over-weight the boundary nodes, and on a hexagonal prism the bounding box is
$4/3$ the crystal.

\item \textbf{The 28-parameter fit is stale} after the atomic-volume correction
of Change~\ref{chg:omega}. The loop densities halve. A refit is planned
separately and is deliberately \emph{not} attempted here.

\item \textbf{DAD is a fitting form, not a generated one.}
\code{dadAnisotropy}/\code{dadZ0} are fitted to reproduce the 0-D symmetric
bias forms rather than derived from $D_c/D_a$, and the diffusion tensors are
isotropic to match ZrMicro. These runs therefore exercise none of the
tensorial-DAD mechanism of the deliverable, and the fit is tied to 573~K.

\item \textbf{SIPN stress weights $w_k(\sigma)$ are not wired.} Exact at zero
applied stress, which is the configured condition for every run reported here.

\item \textbf{\code{Warning: Sum of R2 is not zero}} remains, and is now
\emph{meaningful} rather than merely benign: the reaction map contains
$i{+}3i$ but there is no $4i$ species to receive the product. Those
interstitials are no longer unaccounted --- they are exactly the loop embryos of
Change~\ref{chg:clustering}, and the immobile solve now credits them. ZrMicro
truncates the cluster ladder at the same place for the same reason.

\item \textbf{Near-boundary loop densities are not quantitative.} Cascade
nucleation is spatially uniform, but the only loop-removal channel is
coalescence, driven by the absorbed mobile flux --- which vanishes where
Dirichlet pins the mobile concentrations. The boundary is a sink for mobile
defects but \textbf{not} for loops. Of the 24\,115 CD nodes on the 1~\textmu m
cube, 6522 (27\%) lie on the six Dirichlet faces, and at 21~\dpa\ they carry
\textbf{99.8\%} of the total \aloop\ density. \textbf{Never report a
whole-domain mean of a loop quantity on such a geometry} --- quote the interior
mean, the innermost quartile by distance to the nearest face. The difference is
not cosmetic: at 21~\dpa\ the coupled march's \aloop\ density is $474\times$ the
0-D as a domain mean and $0.75\times$ it in the interior.

That the coupled march reproduces this artifact is a point \emph{in favor of}
the coupling: the same boundary behavior emerges whether MoDELib3's nodal
scheme or ZrMicro's CVODE march advances the immobile population.

\end{enumerate}

% ===========================================================================
\section*{Verification status against the 0-D}
\addcontentsline{toc}{section}{Verification status against the 0-D}
% ===========================================================================

Ratios 3-D / 0-D at interior nodes, from the 31-step run to 30~\dpa. ``Before''
is cascade-only nucleation, ``after'' is with the homogeneous SIA clustering
flux of Change~\ref{chg:clustering}.

\begin{center}
\begin{tabular}{lrrrr}
\toprule
quantity & 1 \dpa & 5 \dpa & 10 \dpa & 30 \dpa \\
\midrule
$C_v$   & 1.073 & 0.982 & 0.994 & 1.004 \\
$C_i$   & 1.145 & 1.112 & 1.133 & 1.148 \\
$C_{2i}$ & 1.012 & 0.412 & 0.394 & 0.388 \\
$C_{3i}$ & 0.977 & 0.400 & 0.383 & 0.377 \\
\cloop\ loop density & 0.999 & 0.999 & 0.999 & 0.999 \\
\aloop\ loop density --- before & 0.426 & 0.547 & 0.540 & 0.533 \\
\aloop\ loop density --- \textbf{after} & \textbf{0.662} & \textbf{0.739} & \textbf{0.718} & \textbf{0.704} \\
\aloop\ loop content --- before & 0.486 & 0.634 & 0.627 & 0.620 \\
\aloop\ loop content --- \textbf{after} & \textbf{0.639} & \textbf{0.811} & \textbf{0.795} & \textbf{0.784} \\
\aloop\ defects/loop --- before & 1.141 & 1.159 & 1.159 & 1.164 \\
\aloop\ defects/loop --- \textbf{after} & \textbf{0.965} & \textbf{1.098} & \textbf{1.106} & \textbf{1.114} \\
\bottomrule
\end{tabular}
\end{center}

\noindent
Interior nodes are selected by the top decile of $C_v$, and \textbf{that
selection must be made ONCE from the mobile field and applied to all
quantities}. Loop density anti-correlates with $C_v$ (corr $= -0.98$), because
near the boundary the mobile species are depleted so loops neither grow nor
coalesce and simply accumulate at their nucleation value. A per-quantity
percentile would pick the boundary for the immobile fields and the interior for
the mobile ones.

$C_v$ agrees to within 2\% beyond 5~\dpa\ and the \cloop\ loop density to 0.1\%
at every dose. $C_i$ runs 11--15\% high throughout, at every dose and in both
runs. The two remaining discrepancies are one story, not two --- see
Part~\ref{part:open}, item~3.

\medskip\noindent
\textbf{Note on scope.} The table above predates the atomic-volume correction
of Change~\ref{chg:omega}. Both codes changed together and by the same factor,
so the \emph{ratios} are expected to survive it, but they have not been
re-measured. They will be, as part of the planned refit.

% ===========================================================================
\section*{Summary of measured efficiency changes}
\addcontentsline{toc}{section}{Summary of measured efficiency changes}
% ===========================================================================

\begin{center}
\begin{tabular}{@{}cp{7.3cm}p{5.4cm}@{}}
\toprule
Change & What & Measured \\
\midrule
\ref{chg:rediscretize} & nodal immobile update, no mass-matrix projection & avoids a CG solve per dose step \\
\ref{chg:anneal-implicit} & vacancy-loop dissolution implicit & unconditional stability (was $\times$35 over) \\
\ref{chg:coal-implicit} & coalescence implicit & unconditional stability (was $\times$3.6 over) \\
\ref{chg:clamp} & clamp deferred out of the Newton loop & $3.5\times$, and quadratic convergence \\
\ref{chg:cvode} & operator split, mobile frozen & $1.0\times10^{10}$ in factorization flops \\
\ref{chg:ad} & AD Jacobian & accuracy, not wall time \\
\ref{chg:quadrature} & accumulators as quadratures & $19\rightarrow13$ standalone; 15 split \\
\ref{chg:workspace} & one SUNDIALS workspace per thread & $\mathbf{1.67\times}$ end to end \\
\ref{chg:elastic} & skip the trivial elastic solve & $\mathbf{1.86\times}$, bit-identical \\
\ref{chg:acir} & remove dead \code{AcIR} assembly & $1.64\times$ \\
\ref{chg:filter} & cached $T$, filtered $A_1$ & $\sim$1.00$\times$ alone \\
\ref{chg:parallel} & parallel element assembly & $1.33\times$ \\
\ref{chg:hoist} & hoist invariant assembly & $1.11\times$ \\
\ref{chg:abcompare} & \ref{chg:acir}--\ref{chg:hoist} on a full march &
  $\mathbf{2.176\times}$ fast, $\mathbf{1.829\times}$ march \\
\bottomrule
\end{tabular}
\end{center}

\noindent
The fast solve remains the bottleneck, and \code{setFromTriplets} is now its
largest single component. The report step is $\sim$47\% of a full run's wall
time, which is the next thing worth attention.

\end{document}
